\noindent Cuatro simulacros originales de \textbf{150 minutos}, con cinco problemas de dos puntos cada uno. Tiempo y puntuación son propuestas de entrenamiento. Las veinte variantes reutilizan familias del banco con datos distintos. Las soluciones están reunidas en la segunda parte, para permitir trabajar sin verlas.
\tableofcontents
\mainmatter
\part{Enunciados}


\chapter{Simulacro 1}

\objetivos{Resolver cinco situaciones de temas diferentes en 150 minutos. Reserva unos 25 minutos por problema y 25 minutos para comprobar.}

Calculadora científica y una hoja de fórmulas propia permitidas. No se necesita ordenador. Cada problema vale 2 puntos: planteamiento e hipótesis, 0,5; ecuaciones y desarrollo, 0,75; resultados con unidades, 0,5; interpretación o comprobación, 0.25. Un error aritmético arrastrado se penaliza una sola vez si el procedimiento posterior es coherente.


\par\medskip\noindent\begin{minipage}{\linewidth}

\section{Problema 1 · 2 puntos}

Una masa se mueve horizontalmente sin tracción y su resistencia es $k v^2$. Deduzca velocidad y recorrido, y calcúlelos al tiempo indicado. Explique si el modelo predice una detención exacta en tiempo finito.

\begin{tabular}{@{}p{.15\linewidth}p{.23\linewidth}p{.54\linewidth}@{}}
\toprule Símbolo & Valor SI & Significado\\\midrule
$m$ & $30\,\mathrm{kg}$ & masa\\
$k$ & $0{,}6\,\mathrm{kg/m}$ & resistencia cuadrática\\
$v_0$ & $10\,\mathrm{m/s}$ & velocidad inicial\\
$t$ & $4\,\mathrm{s}$ & tiempo\\
\bottomrule\end{tabular}

\par\medskip\textbf{Condiciones del modelo.} Velocidad positiva, masa constante, fuerza opuesta al movimiento; no hay otras fuerzas horizontales.

\par\textit{Entrega: ecuaciones justificadas, resultados con unidades y una comprobación física o matemática.}

\end{minipage}\par\medskip

\par\medskip\noindent\begin{minipage}{\linewidth}

\section{Problema 2 · 2 puntos}

Para $\boldsymbol u=(ax,-ay,0)$, determine divergencia, vorticidad, aceleración material y tasa de deformación. Evalúe la aceleración en el punto dado y explique por qué estacionario no significa aceleración nula.

\begin{tabular}{@{}p{.15\linewidth}p{.23\linewidth}p{.54\linewidth}@{}}
\toprule Símbolo & Valor SI & Significado\\\midrule
$a$ & $0{,}8\,\mathrm{s^{-1}}$ & tasa de extensión\\
$x$ & $0{,}25\,\mathrm{m}$ & coordenada x\\
$y$ & $0{,}15\,\mathrm{m}$ & coordenada y\\
\bottomrule\end{tabular}

\par\medskip\textbf{Condiciones del modelo.} Campo cartesiano estacionario, diferenciable en el dominio considerado. No se prescribe aquí una pared ni una geometría de entrada.

\par\textit{Entrega: ecuaciones justificadas, resultados con unidades y una comprobación física o matemática.}

\end{minipage}\par\medskip

\par\medskip\noindent\begin{minipage}{\linewidth}

\section{Problema 3 · 2 puntos}

Una compuerta rectangular está articulada en su borde superior, situado bajo una superficie libre. Determine fuerza hidrostática, punto de aplicación y fuerza normal que debe ejercerse en el extremo inferior para equilibrarla.

\begin{tabular}{@{}p{.15\linewidth}p{.23\linewidth}p{.54\linewidth}@{}}
\toprule Símbolo & Valor SI & Significado\\\midrule
$b$ & $1\,\mathrm{m}$ & anchura\\
$L$ & $1{,}8\,\mathrm{m}$ & longitud sobre el panel\\
$h_0$ & $0{,}3\,\mathrm{m}$ & profundidad de la articulación\\
$\theta$ & $0{,}785398$ & ángulo con la horizontal\\
$\rho$ & $1000\,\mathrm{kg/m^3}$ & densidad\\
\bottomrule\end{tabular}

\par\medskip\textbf{Condiciones del modelo.} Líquido en reposo, presión atmosférica en la cara seca y superficie libre. Se omiten peso del panel y rozamiento de la articulación.

\par\textit{Entrega: ecuaciones justificadas, resultados con unidades y una comprobación física o matemática.}

\end{minipage}\par\medskip

\par\medskip\noindent\begin{minipage}{\linewidth}

\section{Problema 4 · 2 puntos}

Un capilar transporta el caudal dado. Calcule caída de presión, velocidad media, Reynolds, esfuerzo en pared y potencia disipada. Compruebe la coherencia del supuesto laminar.

\begin{tabular}{@{}p{.15\linewidth}p{.23\linewidth}p{.54\linewidth}@{}}
\toprule Símbolo & Valor SI & Significado\\\midrule
$R$ & $0{,}0012\,\mathrm{m}$ & radio\\
$L$ & $1\,\mathrm{m}$ & longitud\\
$\rho$ & $980\,\mathrm{kg/m^3}$ & densidad\\
$\mu$ & $0{,}012\,\mathrm{Pa\,s}$ & viscosidad\\
$Q$ & $1{,}5\times10^{-6}\,\mathrm{m^3/s}$ & caudal\\
\bottomrule\end{tabular}

\par\medskip\textbf{Condiciones del modelo.} Tubo circular horizontal rígido; fluido incompresible newtoniano y flujo estacionario plenamente desarrollado.

\par\textit{Entrega: ecuaciones justificadas, resultados con unidades y una comprobación física o matemática.}

\end{minipage}\par\medskip

\par\medskip\noindent\begin{minipage}{\linewidth}

\section{Problema 5 · 2 puntos}

Con el gradiente horizontal de presión y el parámetro de Coriolis dados, determine ambas componentes geostróficas y el número de Rossby usando la escala horizontal indicada.

\begin{tabular}{@{}p{.15\linewidth}p{.23\linewidth}p{.54\linewidth}@{}}
\toprule Símbolo & Valor SI & Significado\\\midrule
$f$ & $0{,}0001\,\mathrm{s^{-1}}$ & parámetro de Coriolis\\
$\rho$ & $1{,}2\,\mathrm{kg/m^3}$ & densidad\\
$\partial_xp$ & $0{,}0012\,\mathrm{Pa/m}$ & gradiente hacia x\\
$\partial_yp$ & $-0{,}00048\,\mathrm{Pa/m}$ & gradiente hacia y\\
$L$ & $1{,}8\times10^{6}\,\mathrm{m}$ & escala horizontal\\
\bottomrule\end{tabular}

\par\medskip\textbf{Condiciones del modelo.} Plano local con $f>0$ constante, flujo horizontal de gran escala sin fricción ni aceleración en el balance principal. $x,y,z$ forman una base dextrógira.

\par\textit{Entrega: ecuaciones justificadas, resultados con unidades y una comprobación física o matemática.}

\end{minipage}\par\medskip

\resumen{Antes de entregar, comprueba signos de las fuerzas, definición de cada número adimensional y carácter absoluto o manométrico de la presión.}

\chapter{Simulacro 2}

\objetivos{Resolver cinco situaciones de temas diferentes en 150 minutos. Reserva unos 25 minutos por problema y 25 minutos para comprobar.}

Calculadora científica y una hoja de fórmulas propia permitidas. No se necesita ordenador. Cada problema vale 2 puntos: planteamiento e hipótesis, 0,5; ecuaciones y desarrollo, 0,75; resultados con unidades, 0,5; interpretación o comprobación, 0.25. Un error aritmético arrastrado se penaliza una sola vez si el procedimiento posterior es coherente.


\par\medskip\noindent\begin{minipage}{\linewidth}

\section{Problema 1 · 2 puntos}

Una partícula se libera desde el reposo. Tomando positivo hacia abajo, obtenga velocidad límite, tiempo de relajación, velocidad y distancia recorrida. Distinga la energía potencial perdida de la energía cinética adquirida.

\begin{tabular}{@{}p{.15\linewidth}p{.23\linewidth}p{.54\linewidth}@{}}
\toprule Símbolo & Valor SI & Significado\\\midrule
$m$ & $0{,}06\,\mathrm{kg}$ & masa\\
$b$ & $0{,}15\,\mathrm{kg/s}$ & coeficiente lineal\\
$t$ & $0{,}9\,\mathrm{s}$ & tiempo\\
\bottomrule\end{tabular}

\par\medskip\textbf{Condiciones del modelo.} Fuerza resistiva $-bv$, gravedad uniforme y flotación despreciada; el coeficiente se toma constante.

\par\textit{Entrega: ecuaciones justificadas, resultados con unidades y una comprobación física o matemática.}

\end{minipage}\par\medskip

\par\medskip\noindent\begin{minipage}{\linewidth}

\section{Problema 2 · 2 puntos}

Un prisma homogéneo flota con su cara ancha horizontal. Obtenga calado y altura metacéntrica para balanceo alrededor de su eje longitudinal. Estime el momento restaurador para un ángulo pequeño y decida estabilidad inicial.

\begin{tabular}{@{}p{.15\linewidth}p{.23\linewidth}p{.54\linewidth}@{}}
\toprule Símbolo & Valor SI & Significado\\\midrule
$b$ & $0{,}45\,\mathrm{m}$ & manga\\
$h$ & $0{,}22\,\mathrm{m}$ & altura\\
$L$ & $1\,\mathrm{m}$ & longitud\\
$\rho_s$ & $650\,\mathrm{kg/m^3}$ & densidad del sólido\\
$\rho_f$ & $1000\,\mathrm{kg/m^3}$ & densidad del líquido\\
$\theta$ & $0{,}05$ & ángulo en radianes\\
\bottomrule\end{tabular}

\par\medskip\textbf{Condiciones del modelo.} Prisma rígido homogéneo, líquido en reposo y escora infinitesimal; el eje de balanceo es paralelo a la longitud.

\par\textit{Entrega: ecuaciones justificadas, resultados con unidades y una comprobación física o matemática.}

\end{minipage}\par\medskip

\par\medskip\noindent\begin{minipage}{\linewidth}

\section{Problema 3 · 2 puntos}

Una pared plana oscila sinusoidalmente junto a un líquido profundo. Halle espesor de penetración, amplitud y retraso de fase a la distancia dada, esfuerzo máximo en pared y potencia media suministrada por unidad de área.

\begin{tabular}{@{}p{.15\linewidth}p{.23\linewidth}p{.54\linewidth}@{}}
\toprule Símbolo & Valor SI & Significado\\\midrule
$\nu$ & $1\times10^{-5}\,\mathrm{m^2/s}$ & viscosidad cinemática\\
$U$ & $0{,}15\,\mathrm{m/s}$ & amplitud de pared\\
$f$ & $3\,\mathrm{s^{-1}}$ & frecuencia en ciclos por segundo\\
$\rho$ & $900\,\mathrm{kg/m^3}$ & densidad\\
$y$ & $0{,}0008\,\mathrm{m}$ & distancia a pared\\
\bottomrule\end{tabular}

\par\medskip\textbf{Condiciones del modelo.} Fluido newtoniano semiinfinito, flujo paralelo y régimen periódico establecido, sin gradiente de presión. $y$ apunta hacia el líquido.

\par\textit{Entrega: ecuaciones justificadas, resultados con unidades y una comprobación física o matemática.}

\end{minipage}\par\medskip

\par\medskip\noindent\begin{minipage}{\linewidth}

\section{Problema 4 · 2 puntos}

Un codo desvía un chorro un ángulo $\theta$ desde $+x$ hacia $+y$, manteniendo la sección. Calcule la fuerza que el líquido ejerce sobre el codo, indicando componentes y signo.

\begin{tabular}{@{}p{.15\linewidth}p{.23\linewidth}p{.54\linewidth}@{}}
\toprule Símbolo & Valor SI & Significado\\\midrule
$Q$ & $0{,}015\,\mathrm{m^3/s}$ & caudal\\
$D$ & $0{,}035\,\mathrm{m}$ & diámetro\\
$\rho$ & $1000\,\mathrm{kg/m^3}$ & densidad\\
$p$ & $100000\,\mathrm{Pa}$ & presión manométrica de entrada\\
$\theta$ & $1{,}5708$ & ángulo de salida en radianes\\
\bottomrule\end{tabular}

\par\medskip\textbf{Condiciones del modelo.} Balance estacionario con perfiles uniformes, salida atmosférica y peso del volumen de líquido omitido. Las presiones se expresan respecto de la atmósfera.

\par\textit{Entrega: ecuaciones justificadas, resultados con unidades y una comprobación física o matemática.}

\end{minipage}\par\medskip

\par\medskip\noindent\begin{minipage}{\linewidth}

\section{Problema 5 · 2 puntos}

Calcule pérdidas en un tubo con la ley laminar, indicando los factores de Darcy y Fanning. Verifique el resultado mediante Poiseuille y explique el error de mezclar las dos convenciones.

\begin{tabular}{@{}p{.15\linewidth}p{.23\linewidth}p{.54\linewidth}@{}}
\toprule Símbolo & Valor SI & Significado\\\midrule
$\rho$ & $900\,\mathrm{kg/m^3}$ & densidad\\
$\mu$ & $0{,}009\,\mathrm{Pa\,s}$ & viscosidad\\
$D$ & $0{,}012\,\mathrm{m}$ & diámetro\\
$V$ & $0{,}8\,\mathrm{m/s}$ & velocidad media\\
$L$ & $3\,\mathrm{m}$ & longitud\\
\bottomrule\end{tabular}

\par\medskip\textbf{Condiciones del modelo.} Tubo circular horizontal, flujo laminar newtoniano plenamente desarrollado. Se omiten pérdidas locales y de entrada.

\par\textit{Entrega: ecuaciones justificadas, resultados con unidades y una comprobación física o matemática.}

\end{minipage}\par\medskip

\resumen{Antes de entregar, comprueba signos de las fuerzas, definición de cada número adimensional y carácter absoluto o manométrico de la presión.}

\chapter{Simulacro 3}

\objetivos{Resolver cinco situaciones de temas diferentes en 150 minutos. Reserva unos 25 minutos por problema y 25 minutos para comprobar.}

Calculadora científica y una hoja de fórmulas propia permitidas. No se necesita ordenador. Cada problema vale 2 puntos: planteamiento e hipótesis, 0,5; ecuaciones y desarrollo, 0,75; resultados con unidades, 0,5; interpretación o comprobación, 0.25. Un error aritmético arrastrado se penaliza una sola vez si el procedimiento posterior es coherente.


\par\medskip\noindent\begin{minipage}{\linewidth}

\section{Problema 1 · 2 puntos}

Una placa arrastra una capa de líquido. A partir de la fuerza tangencial, obtenga viscosidad dinámica y cinemática, esfuerzo y potencia disipada. Calcule un Reynolds basado en el espesor.

\begin{tabular}{@{}p{.15\linewidth}p{.23\linewidth}p{.54\linewidth}@{}}
\toprule Símbolo & Valor SI & Significado\\\midrule
$F$ & $4{,}2\,\mathrm{N}$ & fuerza tangencial\\
$A$ & $0{,}14\,\mathrm{m^2}$ & área mojada\\
$h$ & $0{,}0015\,\mathrm{m}$ & separación\\
$U$ & $0{,}5\,\mathrm{m/s}$ & velocidad de la placa\\
$\rho$ & $950\,\mathrm{kg/m^3}$ & densidad\\
\bottomrule\end{tabular}

\par\medskip\textbf{Condiciones del modelo.} Fluido newtoniano, perfil lineal, no deslizamiento, placa opuesta fija y gradiente de presión nulo. Fuerza neta de una sola cara mojada.

\par\textit{Entrega: ecuaciones justificadas, resultados con unidades y una comprobación física o matemática.}

\end{minipage}\par\medskip

\par\medskip\noindent\begin{minipage}{\linewidth}

\section{Problema 2 · 2 puntos}

Un avión en vuelo horizontal tiene $C_D=C_{D0}+KC_L^2$. Halle el $C_L$, velocidad y resistencia de máximo $L/D$. Separe la contribución inducida $K_i=1/(\pi eAR)$ de otras contribuciones cuadráticas.

\begin{tabular}{@{}p{.15\linewidth}p{.23\linewidth}p{.54\linewidth}@{}}
\toprule Símbolo & Valor SI & Significado\\\midrule
$C_{D0}$ & $0{,}028$ & coeficiente constante\\
$K$ & $0{,}052$ & coeficiente cuadrático total\\
$e$ & $0{,}9$ & factor inducido dado\\
$AR$ & $9$ & alargamiento\\
$W$ & $5000\,\mathrm{N}$ & peso\\
$\rho$ & $1{,}1\,\mathrm{kg/m^3}$ & densidad\\
$S$ & $14\,\mathrm{m^2}$ & superficie\\
\bottomrule\end{tabular}

\par\medskip\textbf{Condiciones del modelo.} Vuelo estacionario horizontal, $L=W$, polar válida en el intervalo de sustentación empleado; datos sintéticos sin atribución a un avión.

\par\textit{Entrega: ecuaciones justificadas, resultados con unidades y una comprobación física o matemática.}

\end{minipage}\par\medskip

\par\medskip\noindent\begin{minipage}{\linewidth}

\section{Problema 3 · 2 puntos}

En un dominio anular se prescribe $u_r=C/r$, $u_\theta=K/r$. Halle caudal radial por longitud, circulación y aumento de presión entre dos radios. Explique por qué puede existir circulación con rotacional local nulo.

\begin{tabular}{@{}p{.15\linewidth}p{.23\linewidth}p{.54\linewidth}@{}}
\toprule Símbolo & Valor SI & Significado\\\midrule
$C$ & $0{,}004\,\mathrm{m^2/s}$ & intensidad radial\\
$K$ & $0{,}008\,\mathrm{m^2/s}$ & intensidad azimutal\\
$r_1$ & $0{,}12\,\mathrm{m}$ & radio interior\\
$r_2$ & $0{,}24\,\mathrm{m}$ & radio exterior\\
$\rho$ & $1000\,\mathrm{kg/m^3}$ & densidad\\
\bottomrule\end{tabular}

\par\medskip\textbf{Condiciones del modelo.} Flujo ideal estacionario bidimensional en $r_1\le r\le r_2$, sin incluir la singularidad del origen; densidad constante.

\par\textit{Entrega: ecuaciones justificadas, resultados con unidades y una comprobación física o matemática.}

\end{minipage}\par\medskip

\par\medskip\noindent\begin{minipage}{\linewidth}

\section{Problema 4 · 2 puntos}

Un depósito de sección constante descarga por un orificio. Deduzca el tiempo para bajar entre los niveles dados y el caudal inicial, usando el coeficiente de descarga proporcionado.

\begin{tabular}{@{}p{.15\linewidth}p{.23\linewidth}p{.54\linewidth}@{}}
\toprule Símbolo & Valor SI & Significado\\\midrule
$A_t$ & $1\,\mathrm{m^2}$ & sección del depósito\\
$A_o$ & $0{,}0005\,\mathrm{m^2}$ & área del orificio\\
$C_d$ & $0{,}62$ & coeficiente de descarga dado\\
$h_0$ & $2{,}5\,\mathrm{m}$ & nivel inicial sobre orificio\\
$h_1$ & $0{,}8\,\mathrm{m}$ & nivel final sobre orificio\\
\bottomrule\end{tabular}

\par\medskip\textbf{Condiciones del modelo.} Superficie y salida atmosféricas, líquido incompresible, orificio pequeño respecto del depósito y coeficiente constante entre los niveles dados.

\par\textit{Entrega: ecuaciones justificadas, resultados con unidades y una comprobación física o matemática.}

\end{minipage}\par\medskip

\par\medskip\noindent\begin{minipage}{\linewidth}

\section{Problema 5 · 2 puntos}

Para $u=ax+by$, $v=-ay$, $w=0$, calcule tensor de esfuerzos, tracción normal y tangencial sobre el plano de normal $(1,1,0)/\sqrt2$ y disipación viscosa. Use tangente $(-1,1,0)/\sqrt2$.

\begin{tabular}{@{}p{.15\linewidth}p{.23\linewidth}p{.54\linewidth}@{}}
\toprule Símbolo & Valor SI & Significado\\\midrule
$a$ & $1{,}5\,\mathrm{s^{-1}}$ & extensión\\
$b$ & $4\,\mathrm{s^{-1}}$ & cizalla\\
$\mu$ & $0{,}3\,\mathrm{Pa\,s}$ & viscosidad\\
$p$ & $1500\,\mathrm{Pa}$ & presión\\
\bottomrule\end{tabular}

\par\medskip\textbf{Condiciones del modelo.} Fluido newtoniano incompresible isotrópico, esfuerzo $\boldsymbol\sigma=-p\boldsymbol I+2\mu\boldsymbol D$. La tracción se define como $\boldsymbol\sigma\boldsymbol n$.

\par\textit{Entrega: ecuaciones justificadas, resultados con unidades y una comprobación física o matemática.}

\end{minipage}\par\medskip

\resumen{Antes de entregar, comprueba signos de las fuerzas, definición de cada número adimensional y carácter absoluto o manométrico de la presión.}

\chapter{Simulacro 4}

\objetivos{Resolver cinco situaciones de temas diferentes en 150 minutos. Reserva unos 25 minutos por problema y 25 minutos para comprobar.}

Calculadora científica y una hoja de fórmulas propia permitidas. No se necesita ordenador. Cada problema vale 2 puntos: planteamiento e hipótesis, 0,5; ecuaciones y desarrollo, 0,75; resultados con unidades, 0,5; interpretación o comprobación, 0.25. Un error aritmético arrastrado se penaliza una sola vez si el procedimiento posterior es coherente.


\par\medskip\noindent\begin{minipage}{\linewidth}

\section{Problema 1 · 2 puntos}

Una masa unida a un muelle y un amortiguador parte de una elongación sin velocidad. Obtenga su ecuación, determine el régimen y calcule frecuencia amortiguada y amplitud de la envolvente al tiempo dado.

\begin{tabular}{@{}p{.15\linewidth}p{.23\linewidth}p{.54\linewidth}@{}}
\toprule Símbolo & Valor SI & Significado\\\midrule
$m$ & $3\,\mathrm{kg}$ & masa\\
$k$ & $120\,\mathrm{N/m}$ & rigidez\\
$b$ & $6\,\mathrm{kg/s}$ & amortiguamiento\\
$x_0$ & $0{,}04\,\mathrm{m}$ & elongación inicial\\
$t$ & $1\,\mathrm{s}$ & tiempo\\
\bottomrule\end{tabular}

\par\medskip\textbf{Condiciones del modelo.} Un grado de libertad, muelle lineal y fuerza disipativa $-b\dot x$; no hay excitación exterior.

\par\textit{Entrega: ecuaciones justificadas, resultados con unidades y una comprobación física o matemática.}

\end{minipage}\par\medskip

\par\medskip\noindent\begin{minipage}{\linewidth}

\section{Problema 2 · 2 puntos}

Un flujo de baja velocidad está forzado en un tiempo muy corto. Calcule Reynolds, Strouhal y su producto para decidir qué términos inerciales pueden despreciarse. Estime la escala viscosa de presión.

\begin{tabular}{@{}p{.15\linewidth}p{.23\linewidth}p{.54\linewidth}@{}}
\toprule Símbolo & Valor SI & Significado\\\midrule
$\rho$ & $1000\,\mathrm{kg/m^3}$ & densidad\\
$\mu$ & $1\,\mathrm{Pa\,s}$ & viscosidad\\
$L$ & $0{,}0012\,\mathrm{m}$ & longitud\\
$U$ & $0{,}0002\,\mathrm{m/s}$ & velocidad\\
$T$ & $0{,}0002\,\mathrm{s}$ & tiempo impuesto\\
\bottomrule\end{tabular}

\par\medskip\textbf{Condiciones del modelo.} Fluido newtoniano, longitud espacial $L$ y escala temporal externa $T$ independientes; comparación de términos de Navier--Stokes.

\par\textit{Entrega: ecuaciones justificadas, resultados con unidades y una comprobación física o matemática.}

\end{minipage}\par\medskip

\par\medskip\noindent\begin{minipage}{\linewidth}

\section{Problema 3 · 2 puntos}

Un carro se aleja en la dirección de un chorro y desvía el agua noventa grados en su sistema relativo. Halle caudal capturado, fuerza propulsora y potencia mecánica. Compruebe la energía en el sistema del laboratorio.

\begin{tabular}{@{}p{.15\linewidth}p{.23\linewidth}p{.54\linewidth}@{}}
\toprule Símbolo & Valor SI & Significado\\\midrule
$V$ & $24\,\mathrm{m/s}$ & velocidad del chorro en laboratorio\\
$U$ & $6\,\mathrm{m/s}$ & velocidad constante del carro\\
$A$ & $0{,}0008\,\mathrm{m^2}$ & sección del chorro\\
$\rho$ & $1000\,\mathrm{kg/m^3}$ & densidad\\
\bottomrule\end{tabular}

\par\medskip\textbf{Condiciones del modelo.} Chorro libre con $V>U$, sin pérdidas en la desviación relativa; la componente longitudinal relativa de salida es nula. No se cuenta todo el caudal emitido como capturado.

\par\textit{Entrega: ecuaciones justificadas, resultados con unidades y una comprobación física o matemática.}

\end{minipage}\par\medskip

\par\medskip\noindent\begin{minipage}{\linewidth}

\section{Problema 4 · 2 puntos}

Una placa infinita arranca instantáneamente con velocidad $U$. Determine velocidad en el líquido, esfuerzo en pared, transporte volumétrico por anchura y escala de penetración al tiempo dado.

\begin{tabular}{@{}p{.15\linewidth}p{.23\linewidth}p{.54\linewidth}@{}}
\toprule Símbolo & Valor SI & Significado\\\midrule
$U$ & $0{,}3\,\mathrm{m/s}$ & velocidad de placa\\
$\nu$ & $1\times10^{-6}\,\mathrm{m^2/s}$ & viscosidad cinemática\\
$\rho$ & $1000\,\mathrm{kg/m^3}$ & densidad\\
$y$ & $0{,}0012\,\mathrm{m}$ & distancia a placa\\
$t$ & $0{,}8\,\mathrm{s}$ & tiempo desde arranque\\
\bottomrule\end{tabular}

\par\medskip\textbf{Condiciones del modelo.} Líquido newtoniano semiinfinito inicialmente en reposo, flujo paralelo sin gradiente de presión y $t>0$.

\par\textit{Entrega: ecuaciones justificadas, resultados con unidades y una comprobación física o matemática.}

\end{minipage}\par\medskip

\par\medskip\noindent\begin{minipage}{\linewidth}

\section{Problema 5 · 2 puntos}

Con las propiedades académicas dadas, calcule difusividades, Prandtl, Grashof y Rayleigh. Compruebe la pequeña variación relativa de densidad y explique por qué Rayleigh por sí solo no fija un umbral universal de inestabilidad.

\begin{tabular}{@{}p{.15\linewidth}p{.23\linewidth}p{.54\linewidth}@{}}
\toprule Símbolo & Valor SI & Significado\\\midrule
$\rho$ & $1000\,\mathrm{kg/m^3}$ & densidad\\
$\mu$ & $0{,}001\,\mathrm{Pa\,s}$ & viscosidad\\
$k_T$ & $0{,}6\,\mathrm{W/(m\,K)}$ & conductividad\\
$c_p$ & $4200\,\mathrm{J/(kg\,K)}$ & calor específico\\
$\beta$ & $0{,}0002\,\mathrm{K^{-1}}$ & expansión térmica\\
$\Delta T$ & $4\,\mathrm{K}$ & diferencia térmica\\
$L$ & $0{,}012\,\mathrm{m}$ & longitud de referencia\\
\bottomrule\end{tabular}

\par\medskip\textbf{Condiciones del modelo.} Propiedades constantes de referencia, expansión lineal de densidad y gravedad uniforme. El enunciado define escalas, sin especificar una cavidad ni condiciones térmicas de contorno.

\par\textit{Entrega: ecuaciones justificadas, resultados con unidades y una comprobación física o matemática.}

\end{minipage}\par\medskip

\resumen{Antes de entregar, comprueba signos de las fuerzas, definición de cada número adimensional y carácter absoluto o manométrico de la presión.}

\part{Soluciones razonadas}

\chapter{Soluciones del simulacro 1}

\objetivos{Comparar el planteamiento antes de consultar los resultados. La rúbrica valora el modelo, no sólo la cifra final.}

\Needspace{.30\textheight}\section{MF-S1-1. Frenado por resistencia cuadrática}

\begin{studiumworked}{MF-S1-1 · Básico}

\textbf{Enunciado.} Una masa se mueve horizontalmente sin tracción y su resistencia es $k v^2$. Deduzca velocidad y recorrido, y calcúlelos al tiempo indicado. Explique si el modelo predice una detención exacta en tiempo finito.

\begin{longtable}{@{}p{.15\linewidth}p{.23\linewidth}p{.54\linewidth}@{}}
\toprule Símbolo & Valor SI & Significado\\\midrule\endhead
$m$ & $30\,\mathrm{kg}$ & masa\\
$k$ & $0{,}6\,\mathrm{kg/m}$ & resistencia cuadrática\\
$v_0$ & $10\,\mathrm{m/s}$ & velocidad inicial\\
$t$ & $4\,\mathrm{s}$ & tiempo\\
\bottomrule\end{longtable}

\textbf{Incógnitas.} Se calculan $v,\ x,\ \Delta E$ y se interpreta el resultado bajo las hipótesis indicadas.

\textbf{Hipótesis.} Velocidad positiva, masa constante, fuerza opuesta al movimiento; no hay otras fuerzas horizontales.

\textbf{Modelo y ecuaciones.}

\[\fitmath{m\dot v=-kv^2,\quad v(t)=\frac{v_0}{1+kv_0t/m},\quad x(t)=\frac{m}{k}\ln(1+kv_0t/m).}\]

\textbf{Desarrollo.}\begin{enumerate}

\item Separar $dv/v^2=-(k/m)dt$ e integrar con $v(0)=v_0$.

\item Integrar la velocidad entre cero y el tiempo de observación para hallar el desplazamiento.

\item El denominador permanece positivo para $t\ge0$: la velocidad se aproxima a cero sin alcanzarlo en tiempo finito.

\end{enumerate}

\textbf{Cálculo de cada magnitud.}

\begin{enumerate}

\item Obtener $v$ con la relación del modelo y los valores anteriores.

\[\fitmath{v=\frac{v_{0}}{1+\frac{k\cdot v_{0}\cdot t}{m}}}\]

\[\fitmath{v=\frac{(10)}{1+\frac{(0{,}6)\cdot (10)\cdot (4)}{(30)}}\approx 5{,}555556\,\mathrm{m/s}}\]

\item Obtener $x$ con la relación del modelo y los valores anteriores.

\[\fitmath{x=\frac{m}{k}\cdot \ln\left(1+\frac{k\cdot v_{0}\cdot t}{m}\right)}\]

\[\fitmath{x=\frac{(30)}{(0{,}6)}\cdot \ln\left(1+\frac{(0{,}6)\cdot (10)\cdot (4)}{(30)}\right)\approx 29{,}38933\,\mathrm{m}}\]

\item Obtener $\Delta E$ con la relación del modelo y los valores anteriores.

\[\fitmath{\mathrm{loss}=\frac{m\cdot \left({v_{0}}^{2}-{v}^{2}\right)}{2}}\]

\[\fitmath{\mathrm{loss}=\frac{(30)\cdot \left({(10)}^{2}-{(5{,}555555556)}^{2}\right)}{2}\approx 1037{,}037\,\mathrm{J}}\]

\end{enumerate}

\textbf{Resultados.}\begin{itemize}

\item $v=5{,}55556\,\mathrm{m/s}$.

\item $x=29{,}3893\,\mathrm{m}$.

\item $\Delta E=1037{,}04\,\mathrm{J}$.

\end{itemize}

\textbf{Interpretación.} Una ley ideal de resistencia deja una cola de movimiento. La detención de un cuerpo real puede estar gobernada después por fuerzas no incluidas.

\textbf{Comprobación.} Derivar $x(t)$ para recuperar $v(t)$ y comprobar $m\dot v+kv^2=0$. La pérdida cinética debe ser positiva.

\textbf{Límites.} No prolongar la ley cuadrática a velocidades en las que cambie el régimen de resistencia.

\end{studiumworked}

\Needspace{.30\textheight}\section{MF-S1-2. Punto de remanso plano}

\begin{studiumworked}{MF-S1-2 · Intermedio}

\textbf{Enunciado.} Para $\boldsymbol u=(ax,-ay,0)$, determine divergencia, vorticidad, aceleración material y tasa de deformación. Evalúe la aceleración en el punto dado y explique por qué estacionario no significa aceleración nula.

\begin{longtable}{@{}p{.15\linewidth}p{.23\linewidth}p{.54\linewidth}@{}}
\toprule Símbolo & Valor SI & Significado\\\midrule\endhead
$a$ & $0{,}8\,\mathrm{s^{-1}}$ & tasa de extensión\\
$x$ & $0{,}25\,\mathrm{m}$ & coordenada x\\
$y$ & $0{,}15\,\mathrm{m}$ & coordenada y\\
\bottomrule\end{longtable}

\textbf{Incógnitas.} Se calculan $u,\ v,\ a_x,\ a_y$ y se interpreta el resultado bajo las hipótesis indicadas.

\textbf{Hipótesis.} Campo cartesiano estacionario, diferenciable en el dominio considerado. No se prescribe aquí una pared ni una geometría de entrada.

\textbf{Modelo y ecuaciones.}

\[\fitmath{\nabla\cdot\boldsymbol u=0,\quad\nabla\times\boldsymbol u=0,\quad \frac{D\boldsymbol u}{Dt}=(a^2x,a^2y,0),\quad \boldsymbol D=\operatorname{diag}(a,-a,0).}\]

\textbf{Desarrollo.}\begin{enumerate}

\item Calcular las derivadas espaciales de cada componente, conservando los signos del campo.

\item Formar $(\boldsymbol u\cdot\nabla)\boldsymbol u$: el signo negativo de $v$ se multiplica por el de $\partial v/\partial y$.

\item La parte simétrica del gradiente mide extensión y compresión; la antisimétrica se anula.

\end{enumerate}

\textbf{Cálculo de cada magnitud.}

\begin{enumerate}

\item Obtener $u$ con la relación del modelo y los valores anteriores.

\[\fitmath{u=a\cdot x}\]

\[\fitmath{u=(0{,}8)\cdot (0{,}25)\approx 0{,}2\,\mathrm{m/s}}\]

\item Obtener $v$ con la relación del modelo y los valores anteriores.

\[\fitmath{v=-a\cdot y}\]

\[\fitmath{v=-(0{,}8)\cdot (0{,}15)\approx -0{,}12\,\mathrm{m/s}}\]

\item Obtener $a_x$ con la relación del modelo y los valores anteriores.

\[\fitmath{\mathrm{ax}={a}^{2}\cdot x}\]

\[\fitmath{\mathrm{ax}={(0{,}8)}^{2}\cdot (0{,}25)\approx 0{,}16\,\mathrm{m/s^{2}}}\]

\item Obtener $a_y$ con la relación del modelo y los valores anteriores.

\[\fitmath{\mathrm{ay}={a}^{2}\cdot y}\]

\[\fitmath{\mathrm{ay}={(0{,}8)}^{2}\cdot (0{,}15)\approx 0{,}096\,\mathrm{m/s^{2}}}\]

\end{enumerate}

\textbf{Resultados.}\begin{itemize}

\item $u=0{,}2\,\mathrm{m/s}$.

\item $v=-0{,}12\,\mathrm{m/s}$.

\item $a_x=0{,}16\,\mathrm{m/s^2}$.

\item $a_y=0{,}096\,\mathrm{m/s^2}$.

\end{itemize}

\textbf{Interpretación.} La partícula acelera por moverse hacia regiones con otra velocidad, aunque el observador euleriano no vea cambios temporales.

\textbf{Comprobación.} Integrar las trayectorias $x=x_0e^{at}$, $y=y_0e^{-at}$ y comprobar su segunda derivada. El producto $xy$ permanece constante.

\textbf{Límites.} El campo ideal no describe por sí solo el flujo completo alrededor de un cuerpo con no deslizamiento.

\end{studiumworked}

\Needspace{.30\textheight}\section{MF-S1-3. Compuerta inclinada articulada}

\begin{studiumworked}{MF-S1-3 · Intermedio}

\textbf{Enunciado.} Una compuerta rectangular está articulada en su borde superior, situado bajo una superficie libre. Determine fuerza hidrostática, punto de aplicación y fuerza normal que debe ejercerse en el extremo inferior para equilibrarla.

\begin{longtable}{@{}p{.15\linewidth}p{.23\linewidth}p{.54\linewidth}@{}}
\toprule Símbolo & Valor SI & Significado\\\midrule\endhead
$b$ & $1\,\mathrm{m}$ & anchura\\
$L$ & $1{,}8\,\mathrm{m}$ & longitud sobre el panel\\
$h_0$ & $0{,}3\,\mathrm{m}$ & profundidad de la articulación\\
$\theta$ & $0{,}785398$ & ángulo con la horizontal\\
$\rho$ & $1000\,\mathrm{kg/m^3}$ & densidad\\
\bottomrule\end{longtable}

\textbf{Incógnitas.} Se calculan $h_c,\ F,\ M,\ s_R,\ F_{ap}$ y se interpreta el resultado bajo las hipótesis indicadas.

\textbf{Hipótesis.} Líquido en reposo, presión atmosférica en la cara seca y superficie libre. Se omiten peso del panel y rozamiento de la articulación.

\textbf{Modelo y ecuaciones.}

\[\fitmath{p(s)=\rho g(h_0+s\sin\theta),\quad F=b\int_0^L p(s)ds,\quad M=b\int_0^L sp(s)ds,\quad s_R=M/F.}\]

\textbf{Desarrollo.}\begin{enumerate}

\item Usar presión manométrica: la atmosférica se cancela entre ambas caras.

\item Integrar por separado fuerza y momento respecto de la bisagra; $s$ mide distancia sobre el panel, no profundidad vertical.

\item Equilibrar momentos: $F_{ap}L=M$. La reacción de la bisagra completa el equilibrio de fuerzas.

\end{enumerate}

\textbf{Cálculo de cada magnitud.}

\begin{enumerate}

\item Obtener $h_c$ con la relación del modelo y los valores anteriores.

\[\fitmath{\mathrm{hc}=h_{0}+\frac{L\cdot \sin\left(\theta\right)}{2}}\]

\[\fitmath{\mathrm{hc}=(0{,}3)+\frac{(1{,}8)\cdot \sin\left((0{,}7853981634)\right)}{2}\approx 0{,}9363961\,\mathrm{m}}\]

\item Obtener $F$ con la relación del modelo y los valores anteriores.

\[\fitmath{F=\rho\cdot g\cdot b\cdot L\cdot \mathrm{hc}}\]

\[\fitmath{F=(1000)\cdot (9{,}81)\cdot (1)\cdot (1{,}8)\cdot (0{,}9363961031)\approx 16534{,}88\,\mathrm{N}}\]

\item Obtener $M$ con la relación del modelo y los valores anteriores.

\[\fitmath{M=\rho\cdot g\cdot b\cdot \left(\frac{h_{0}\cdot {L}^{2}}{2}+\frac{\sin\left(\theta\right)\cdot {L}^{3}}{3}\right)}\]

\[\fitmath{M=(1000)\cdot (9{,}81)\cdot (1)\cdot \left(\frac{(0{,}3)\cdot {(1{,}8)}^{2}}{2}+\frac{\sin\left((0{,}7853981634)\right)\cdot {(1{,}8)}^{3}}{3}\right)\approx 18252{,}64\,\mathrm{N\,m}}\]

\item Obtener $s_R$ con la relación del modelo y los valores anteriores.

\[\fitmath{\mathrm{sr}=\frac{M}{F}}\]

\[\fitmath{\mathrm{sr}=\frac{(18252{,}63887)}{(16534{,}88239)}\approx 1{,}103887\,\mathrm{m}}\]

\item Obtener $F_{ap}$ con la relación del modelo y los valores anteriores.

\[\fitmath{\mathrm{hold}=\frac{M}{L}}\]

\[\fitmath{\mathrm{hold}=\frac{(18252{,}63887)}{(1{,}8)}\approx 10140{,}35\,\mathrm{N}}\]

\end{enumerate}

\textbf{Resultados.}\begin{itemize}

\item $h_c=0{,}936396\,\mathrm{m}$.

\item $F=16534{,}9\,\mathrm{N}$.

\item $M=18252{,}6\,\mathrm{N\,m}$.

\item $s_R=1{,}10389\,\mathrm{m}$.

\item $F_{ap}=10140{,}4\,\mathrm{N}$.

\end{itemize}

\textbf{Interpretación.} El centro de presiones queda por debajo del centro geométrico porque la presión crece con la profundidad.

\textbf{Comprobación.} Comprobar $L/2<s_R<L$ y recuperar $s_R=2L/3$ cuando $h_0=0$. Una cuadratura de $p(s)$ reproduce fuerza y momento.

\textbf{Límites.} La carga de equilibrio es normal al panel y se aplica en su extremo; otra dirección cambia el brazo de palanca.

\end{studiumworked}

\Needspace{.30\textheight}\section{MF-S1-4. Viscosímetro capilar e inversión del caudal}

\begin{studiumworked}{MF-S1-4 · Básico}

\textbf{Enunciado.} Un capilar transporta el caudal dado. Calcule caída de presión, velocidad media, Reynolds, esfuerzo en pared y potencia disipada. Compruebe la coherencia del supuesto laminar.

\begin{longtable}{@{}p{.15\linewidth}p{.23\linewidth}p{.54\linewidth}@{}}
\toprule Símbolo & Valor SI & Significado\\\midrule\endhead
$R$ & $0{,}0012\,\mathrm{m}$ & radio\\
$L$ & $1\,\mathrm{m}$ & longitud\\
$\rho$ & $980\,\mathrm{kg/m^3}$ & densidad\\
$\mu$ & $0{,}012\,\mathrm{Pa\,s}$ & viscosidad\\
$Q$ & $1{,}5\times10^{-6}\,\mathrm{m^3/s}$ & caudal\\
\bottomrule\end{longtable}

\textbf{Incógnitas.} Se calculan $\Delta p,\ \bar u,\ \mathrm{Re}_D,\ \tau_w,\ P$ y se interpreta el resultado bajo las hipótesis indicadas.

\textbf{Hipótesis.} Tubo circular horizontal rígido; fluido incompresible newtoniano y flujo estacionario plenamente desarrollado.

\textbf{Modelo y ecuaciones.}

\[\fitmath{\Delta p=8\mu LQ/(\pi R^4),\quad \bar u=Q/(\pi R^2),\quad \mathrm{Re}_D=2\rho\bar uR/\mu,\quad \tau_w=\Delta pR/(2L).}\]

\textbf{Desarrollo.}\begin{enumerate}

\item Invertir la ley de caudal de Poiseuille; conservar el radio elevado a cuatro.

\item Usar diámetro $2R$ en Reynolds y señalarlo explícitamente.

\item La potencia $\Delta pQ$ se convierte en disipación en el tramo desarrollado; no incluye una bomba real ni sus pérdidas.

\end{enumerate}

\textbf{Cálculo de cada magnitud.}

\begin{enumerate}

\item Obtener $\Delta p$ con la relación del modelo y los valores anteriores.

\[\fitmath{\mathrm{dp}=\frac{8\cdot \mu\cdot L\cdot Q}{\pi\cdot {R}^{4}}}\]

\[\fitmath{\mathrm{dp}=\frac{8\cdot (0{,}012)\cdot (1)\cdot (1{,}5\times10^{-6})}{\pi\cdot {(0{,}0012)}^{4}}\approx 22104{,}85\,\mathrm{Pa}}\]

\item Obtener $\bar u$ con la relación del modelo y los valores anteriores.

\[\fitmath{\mathrm{um}=\frac{Q}{\pi\cdot {R}^{2}}}\]

\[\fitmath{\mathrm{um}=\frac{(1{,}5\times10^{-6})}{\pi\cdot {(0{,}0012)}^{2}}\approx 0{,}3315728\,\mathrm{m/s}}\]

\item Obtener $\mathrm{Re}_D$ con la relación del modelo y los valores anteriores.

\[\fitmath{\mathrm{Re}=\frac{2\cdot \rho\cdot \mathrm{um}\cdot R}{\mu}}\]

\[\fitmath{\mathrm{Re}=\frac{2\cdot (980)\cdot (0{,}3315727981)\cdot (0{,}0012)}{(0{,}012)}\approx 64{,}98827\,}\]

\item Obtener $\tau_w$ con la relación del modelo y los valores anteriores.

\[\fitmath{\tau=\frac{\mathrm{dp}\cdot R}{2\cdot L}}\]

\[\fitmath{\tau=\frac{(22104{,}85321)\cdot (0{,}0012)}{2\cdot (1)}\approx 13{,}26291\,\mathrm{Pa}}\]

\item Obtener $P$ con la relación del modelo y los valores anteriores.

\[\fitmath{\mathrm{power}=\mathrm{dp}\cdot Q}\]

\[\fitmath{\mathrm{power}=(22104{,}85321)\cdot (1{,}5\times10^{-6})\approx 0{,}03315728\,\mathrm{W}}\]

\end{enumerate}

\textbf{Resultados.}\begin{itemize}

\item $\Delta p=22104{,}9\,\mathrm{Pa}$.

\item $\bar u=0{,}331573\,\mathrm{m/s}$.

\item $\mathrm{Re}_D=64{,}9883$.

\item $\tau_w=13{,}2629\,\mathrm{Pa}$.

\item $P=0{,}0331573\,\mathrm{W}$.

\end{itemize}

\textbf{Interpretación.} El radio controla fuertemente la resistencia. La comprobación de Reynolds apoya el régimen supuesto, pero no sustituye la comprobación de longitud de entrada.

\textbf{Comprobación.} Integrar $u(r)=\Delta p(R^2-r^2)/(4\mu L)$ sobre el área reproduce $Q$. El balance axial da la misma tensión en pared.

\textbf{Límites.} No se incorporan correcciones de entrada, deslizamiento, dependencia térmica ni errores de medida del radio.

\end{studiumworked}

\Needspace{.30\textheight}\section{MF-S1-5. Viento geostrófico y signo del gradiente}

\begin{studiumworked}{MF-S1-5 · Intermedio}

\textbf{Enunciado.} Con el gradiente horizontal de presión y el parámetro de Coriolis dados, determine ambas componentes geostróficas y el número de Rossby usando la escala horizontal indicada.

\begin{longtable}{@{}p{.15\linewidth}p{.23\linewidth}p{.54\linewidth}@{}}
\toprule Símbolo & Valor SI & Significado\\\midrule\endhead
$f$ & $0{,}0001\,\mathrm{s^{-1}}$ & parámetro de Coriolis\\
$\rho$ & $1{,}2\,\mathrm{kg/m^3}$ & densidad\\
$\partial_xp$ & $0{,}0012\,\mathrm{Pa/m}$ & gradiente hacia x\\
$\partial_yp$ & $-0{,}00048\,\mathrm{Pa/m}$ & gradiente hacia y\\
$L$ & $1{,}8\times10^{6}\,\mathrm{m}$ & escala horizontal\\
\bottomrule\end{longtable}

\textbf{Incógnitas.} Se calculan $u_g,\ v_g,\ |\boldsymbol u_g|,\ \mathrm{Ro}$ y se interpreta el resultado bajo las hipótesis indicadas.

\textbf{Hipótesis.} Plano local con $f>0$ constante, flujo horizontal de gran escala sin fricción ni aceleración en el balance principal. $x,y,z$ forman una base dextrógira.

\textbf{Modelo y ecuaciones.}

\[\fitmath{f\boldsymbol e_z\times\boldsymbol u_g=-\nabla_hp/\rho,\quad u_g=-\partial_yp/(\rho f),\quad v_g=\partial_xp/(\rho f).}\]

\textbf{Desarrollo.}\begin{enumerate}

\item Expandir el producto vectorial: sus componentes son $(-fv_g,fu_g)$.

\item Resolver cada ecuación conservando el signo del gradiente suministrado.

\item Calcular Rossby con el módulo de la velocidad resultante y valorar la coherencia del balance dominante.

\end{enumerate}

\textbf{Cálculo de cada magnitud.}

\begin{enumerate}

\item Obtener $u_g$ con la relación del modelo y los valores anteriores.

\[\fitmath{\mathrm{ug}=\frac{-\mathrm{py}}{\rho\cdot f}}\]

\[\fitmath{\mathrm{ug}=\frac{-(-0{,}00048)}{(1{,}2)\cdot (0{,}0001)}\approx 4\,\mathrm{m/s}}\]

\item Obtener $v_g$ con la relación del modelo y los valores anteriores.

\[\fitmath{\mathrm{vg}=\frac{\mathrm{px}}{\rho\cdot f}}\]

\[\fitmath{\mathrm{vg}=\frac{(0{,}0012)}{(1{,}2)\cdot (0{,}0001)}\approx 10\,\mathrm{m/s}}\]

\item Obtener $|\boldsymbol u_g|$ con la relación del modelo y los valores anteriores.

\[\fitmath{\mathrm{speed}=\sqrt{{\mathrm{ug}}^{2}+{\mathrm{vg}}^{2}}}\]

\[\fitmath{\mathrm{speed}=\sqrt{{(4)}^{2}+{(10)}^{2}}\approx 10{,}77033\,\mathrm{m/s}}\]

\item Obtener $\mathrm{Ro}$ con la relación del modelo y los valores anteriores.

\[\fitmath{\mathrm{Ro}=\frac{\mathrm{speed}}{f\cdot L}}\]

\[\fitmath{\mathrm{Ro}=\frac{(10{,}77032961)}{(0{,}0001)\cdot (1800000)}\approx 0{,}05983516\,}\]

\end{enumerate}

\textbf{Resultados.}\begin{itemize}

\item $u_g=4\,\mathrm{m/s}$.

\item $v_g=10\,\mathrm{m/s}$.

\item $|\boldsymbol u_g|=10{,}7703\,\mathrm{m/s}$.

\item $\mathrm{Ro}=0{,}0598352$.

\end{itemize}

\textbf{Interpretación.} La velocidad queda perpendicular al gradiente de presión. Rossby pequeño apoya el balance, sin hacerlo exacto ni garantizarlo cerca de fronteras.

\textbf{Comprobación.} Comprobar $\boldsymbol u_g\cdot\nabla_hp=0$ y ambas componentes de la ecuación vectorial.

\textbf{Límites.} La aproximación falla cerca de $f=0$ y necesita términos adicionales para aceleración, curvatura o fricción relevantes.

\end{studiumworked}

\studiumchaptercheck

\chapter{Soluciones del simulacro 2}

\objetivos{Comparar el planteamiento antes de consultar los resultados. La rúbrica valora el modelo, no sólo la cifra final.}

\Needspace{.30\textheight}\section{MF-S2-1. Caída con resistencia lineal}

\begin{studiumworked}{MF-S2-1 · Intermedio}

\textbf{Enunciado.} Una partícula se libera desde el reposo. Tomando positivo hacia abajo, obtenga velocidad límite, tiempo de relajación, velocidad y distancia recorrida. Distinga la energía potencial perdida de la energía cinética adquirida.

\begin{longtable}{@{}p{.15\linewidth}p{.23\linewidth}p{.54\linewidth}@{}}
\toprule Símbolo & Valor SI & Significado\\\midrule\endhead
$m$ & $0{,}06\,\mathrm{kg}$ & masa\\
$b$ & $0{,}15\,\mathrm{kg/s}$ & coeficiente lineal\\
$t$ & $0{,}9\,\mathrm{s}$ & tiempo\\
\bottomrule\end{longtable}

\textbf{Incógnitas.} Se calculan $\tau,\ v_T,\ v,\ z,\ E_d$ y se interpreta el resultado bajo las hipótesis indicadas.

\textbf{Hipótesis.} Fuerza resistiva $-bv$, gravedad uniforme y flotación despreciada; el coeficiente se toma constante.

\textbf{Modelo y ecuaciones.}

\[\fitmath{\tau=m/b,\quad v_T=mg/b,\quad v=v_T(1-e^{-t/\tau}),\quad z=v_T[t-\tau(1-e^{-t/\tau})].}\]

\textbf{Desarrollo.}\begin{enumerate}

\item Resolver $m\dot v=mg-bv$ con factor integrante y condición inicial nula.

\item Integrar la velocidad para obtener el descenso.

\item Aplicar energía mecánica: $mgz-mv^2/2$ es la energía disipada por la resistencia.

\end{enumerate}

\textbf{Cálculo de cada magnitud.}

\begin{enumerate}

\item Obtener $\tau$ con la relación del modelo y los valores anteriores.

\[\fitmath{\tau=\frac{m}{b}}\]

\[\fitmath{\tau=\frac{(0{,}06)}{(0{,}15)}\approx 0{,}4\,\mathrm{s}}\]

\item Obtener $v_T$ con la relación del modelo y los valores anteriores.

\[\fitmath{\mathrm{vt}=\frac{m\cdot g}{b}}\]

\[\fitmath{\mathrm{vt}=\frac{(0{,}06)\cdot (9{,}81)}{(0{,}15)}\approx 3{,}924\,\mathrm{m/s}}\]

\item Obtener $v$ con la relación del modelo y los valores anteriores.

\[\fitmath{v=\mathrm{vt}\cdot \left(1-\exp\left(\frac{-t}{\tau}\right)\right)}\]

\[\fitmath{v=(3{,}924)\cdot \left(1-\exp\left(\frac{-(0{,}9)}{(0{,}4)}\right)\right)\approx 3{,}510413\,\mathrm{m/s}}\]

\item Obtener $z$ con la relación del modelo y los valores anteriores.

\[\fitmath{z=\mathrm{vt}\cdot \left(t-\tau\cdot \left(1-\exp\left(\frac{-t}{\tau}\right)\right)\right)}\]

\[\fitmath{z=(3{,}924)\cdot \left((0{,}9)-(0{,}4)\cdot \left(1-\exp\left(\frac{-(0{,}9)}{(0{,}4)}\right)\right)\right)\approx 2{,}127435\,\mathrm{m}}\]

\item Obtener $E_d$ con la relación del modelo y los valores anteriores.

\[\fitmath{\mathrm{Ed}=m\cdot g\cdot z-\frac{m\cdot {v}^{2}}{2}}\]

\[\fitmath{\mathrm{Ed}=(0{,}06)\cdot (9{,}81)\cdot (2{,}127434623)-\frac{(0{,}06)\cdot {(3{,}510413443)}^{2}}{2}\approx 0{,}8825179\,\mathrm{J}}\]

\end{enumerate}

\textbf{Resultados.}\begin{itemize}

\item $\tau=0{,}4\,\mathrm{s}$.

\item $v_T=3{,}924\,\mathrm{m/s}$.

\item $v=3{,}51041\,\mathrm{m/s}$.

\item $z=2{,}12743\,\mathrm{m}$.

\item $E_d=0{,}882518\,\mathrm{J}$.

\end{itemize}

\textbf{Interpretación.} La fuerza resultante tiende a cero, mientras la potencia gravitatoria termina disipándose por completo.

\textbf{Comprobación.} La solución satisface $v(0)=z(0)=0$ y $\dot v(0)=g$. Para $t\gg\tau$, $v\to v_T$; comprobar $E_d\ge0$.

\textbf{Límites.} La ley lineal es una hipótesis del enunciado; no se atribuye a cualquier esfera o número de Reynolds.

\end{studiumworked}

\Needspace{.30\textheight}\section{MF-S2-2. Estabilidad de un prisma flotante}

\begin{studiumworked}{MF-S2-2 · Avanzado}

\textbf{Enunciado.} Un prisma homogéneo flota con su cara ancha horizontal. Obtenga calado y altura metacéntrica para balanceo alrededor de su eje longitudinal. Estime el momento restaurador para un ángulo pequeño y decida estabilidad inicial.

\begin{longtable}{@{}p{.15\linewidth}p{.23\linewidth}p{.54\linewidth}@{}}
\toprule Símbolo & Valor SI & Significado\\\midrule\endhead
$b$ & $0{,}45\,\mathrm{m}$ & manga\\
$h$ & $0{,}22\,\mathrm{m}$ & altura\\
$L$ & $1\,\mathrm{m}$ & longitud\\
$\rho_s$ & $650\,\mathrm{kg/m^3}$ & densidad del sólido\\
$\rho_f$ & $1000\,\mathrm{kg/m^3}$ & densidad del líquido\\
$\theta$ & $0{,}05$ & ángulo en radianes\\
\bottomrule\end{longtable}

\textbf{Incógnitas.} Se calculan $d,\ BM,\ BG,\ GM,\ W,\ M_R$ y se interpreta el resultado bajo las hipótesis indicadas.

\textbf{Hipótesis.} Prisma rígido homogéneo, líquido en reposo y escora infinitesimal; el eje de balanceo es paralelo a la longitud.

\textbf{Modelo y ecuaciones.}

\[\fitmath{d=\rho_sh/\rho_f,\quad BM=I_{wp}/\mathcal V=b^2/(12d),\quad BG=(h-d)/2,\quad GM=BM-BG,\quad M_R\simeq-WGM\theta.}\]

\textbf{Desarrollo.}\begin{enumerate}

\item Igualar peso y empuje para obtener volumen desplazado y calado.

\item Usar $I_{wp}=Lb^3/12$ para el eje longitudinal; no intercambiar manga y longitud.

\item Calcular $GM$ y conservar el signo del momento: un $GM$ positivo produce un par opuesto a la escora.

\end{enumerate}

\textbf{Cálculo de cada magnitud.}

\begin{enumerate}

\item Obtener $d$ con la relación del modelo y los valores anteriores.

\[\fitmath{\mathrm{draft}=\frac{\mathrm{rs}}{\mathrm{rf}}\cdot h}\]

\[\fitmath{\mathrm{draft}=\frac{(650)}{(1000)}\cdot (0{,}22)\approx 0{,}143\,\mathrm{m}}\]

\item Obtener $BM$ con la relación del modelo y los valores anteriores.

\[\fitmath{\mathrm{BM}=\frac{{b}^{2}}{12\cdot \mathrm{draft}}}\]

\[\fitmath{\mathrm{BM}=\frac{{(0{,}45)}^{2}}{12\cdot (0{,}143)}\approx 0{,}118007\,\mathrm{m}}\]

\item Obtener $BG$ con la relación del modelo y los valores anteriores.

\[\fitmath{\mathrm{BG}=\frac{h-\mathrm{draft}}{2}}\]

\[\fitmath{\mathrm{BG}=\frac{(0{,}22)-(0{,}143)}{2}\approx 0{,}0385\,\mathrm{m}}\]

\item Obtener $GM$ con la relación del modelo y los valores anteriores.

\[\fitmath{\mathrm{GM}=\mathrm{BM}-\mathrm{BG}}\]

\[\fitmath{\mathrm{GM}=(0{,}118006993)-(0{,}0385)\approx 0{,}07950699\,\mathrm{m}}\]

\item Obtener $W$ con la relación del modelo y los valores anteriores.

\[\fitmath{W=\mathrm{rs}\cdot L\cdot b\cdot h\cdot g}\]

\[\fitmath{W=(650)\cdot (1)\cdot (0{,}45)\cdot (0{,}22)\cdot (9{,}81)\approx 631{,}2735\,\mathrm{N}}\]

\item Obtener $M_R$ con la relación del modelo y los valores anteriores.

\[\fitmath{\mathrm{moment}=-W\cdot \mathrm{GM}\cdot \theta}\]

\[\fitmath{\mathrm{moment}=-(631{,}2735)\cdot (0{,}07950699301)\cdot (0{,}05)\approx -2{,}509533\,\mathrm{N\,m}}\]

\end{enumerate}

\textbf{Resultados.}\begin{itemize}

\item $d=0{,}143\,\mathrm{m}$.

\item $BM=0{,}118007\,\mathrm{m}$.

\item $BG=0{,}0385\,\mathrm{m}$.

\item $GM=0{,}079507\,\mathrm{m}$.

\item $W=631{,}274\,\mathrm{N}$.

\item $M_R=-2{,}50953\,\mathrm{N\,m}$.

\end{itemize}

\textbf{Interpretación.} La condición de flotación no garantiza estabilidad. En este caso la altura metacéntrica positiva da estabilidad inicial de balanceo.

\textbf{Comprobación.} El empuje calculado con el calado debe igualar $W$. Para $GM=0$ desaparece el término restaurador lineal, sin decidir la estabilidad de orden superior.

\textbf{Límites.} No predice estabilidad a grandes ángulos, inundación, balanceo con oleaje ni estabilidad alrededor de otro eje.

\end{studiumworked}

\Needspace{.30\textheight}\section{MF-S2-3. Pared oscilante: segundo problema de Stokes}

\begin{studiumworked}{MF-S2-3 · Intermedio}

\textbf{Enunciado.} Una pared plana oscila sinusoidalmente junto a un líquido profundo. Halle espesor de penetración, amplitud y retraso de fase a la distancia dada, esfuerzo máximo en pared y potencia media suministrada por unidad de área.

\begin{longtable}{@{}p{.15\linewidth}p{.23\linewidth}p{.54\linewidth}@{}}
\toprule Símbolo & Valor SI & Significado\\\midrule\endhead
$\nu$ & $1\times10^{-5}\,\mathrm{m^2/s}$ & viscosidad cinemática\\
$U$ & $0{,}15\,\mathrm{m/s}$ & amplitud de pared\\
$f$ & $3\,\mathrm{s^{-1}}$ & frecuencia en ciclos por segundo\\
$\rho$ & $900\,\mathrm{kg/m^3}$ & densidad\\
$y$ & $0{,}0008\,\mathrm{m}$ & distancia a pared\\
\bottomrule\end{longtable}

\textbf{Incógnitas.} Se calculan $\omega,\ \delta,\ |u|,\ \varphi,\ |\tau_w|,\ \overline P/A$ y se interpreta el resultado bajo las hipótesis indicadas.

\textbf{Hipótesis.} Fluido newtoniano semiinfinito, flujo paralelo y régimen periódico establecido, sin gradiente de presión. $y$ apunta hacia el líquido.

\textbf{Modelo y ecuaciones.}

\[\fitmath{\delta=\sqrt{2\nu/\omega},\quad u=Ue^{-y/\delta}\cos(\omega t-y/\delta),\quad |\tau_w|=\mu U\sqrt{\omega/\nu},\quad \overline P/A=\mu U^2/(2\delta).}\]

\textbf{Desarrollo.}\begin{enumerate}

\item Convertir ciclos por segundo a frecuencia angular $\omega=2\pi f$.

\item La solución compleja tiene un factor de amortiguación y otro de fase con la misma escala espacial.

\item Derivar en pared. La amplitud de esfuerzo contiene $\sqrt2/\delta$; al promediar la potencia se conserva sólo la parte en fase con la velocidad.

\end{enumerate}

\textbf{Cálculo de cada magnitud.}

\begin{enumerate}

\item Obtener $\omega$ con la relación del modelo y los valores anteriores.

\[\fitmath{\omega=2\cdot \pi\cdot f}\]

\[\fitmath{\omega=2\cdot \pi\cdot (3)\approx 18{,}84956\,\mathrm{1/s}}\]

\item Obtener $\delta$ con la relación del modelo y los valores anteriores.

\[\fitmath{\delta=\sqrt{\frac{2\cdot \nu}{\omega}}}\]

\[\fitmath{\delta=\sqrt{\frac{2\cdot (1\times10^{-5})}{(18{,}84955592)}}\approx 0{,}001030065\,\mathrm{m}}\]

\item Obtener $|u|$ con la relación del modelo y los valores anteriores.

\[\fitmath{\mathrm{amp}=U\cdot \exp\left(\frac{-y}{\delta}\right)}\]

\[\fitmath{\mathrm{amp}=(0{,}15)\cdot \exp\left(\frac{-(0{,}0008)}{(0{,}001030064539)}\right)\approx 0{,}06899161\,\mathrm{m/s}}\]

\item Obtener $\varphi$ con la relación del modelo y los valores anteriores.

\[\fitmath{\mathrm{phase}=\frac{y}{\delta}}\]

\[\fitmath{\mathrm{phase}=\frac{(0{,}0008)}{(0{,}001030064539)}\approx 0{,}7766504\,}\]

\item Obtener $|\tau_w|$ con la relación del modelo y los valores anteriores.

\[\fitmath{\tau=\rho\cdot \nu\cdot U\cdot \sqrt{\frac{\omega}{\nu}}}\]

\[\fitmath{\tau=(900)\cdot (1\times10^{-5})\cdot (0{,}15)\cdot \sqrt{\frac{(18{,}84955592)}{(1\times10^{-5})}}\approx 1{,}853465\,\mathrm{Pa}}\]

\item Obtener $\overline P/A$ con la relación del modelo y los valores anteriores.

\[\fitmath{\mathrm{power}=\frac{\rho\cdot \nu\cdot {U}^{2}}{2\cdot \delta}}\]

\[\fitmath{\mathrm{power}=\frac{(900)\cdot (1\times10^{-5})\cdot {(0{,}15)}^{2}}{2\cdot (0{,}001030064539)}\approx 0{,}09829481\,\mathrm{W/m^{2}}}\]

\end{enumerate}

\textbf{Resultados.}\begin{itemize}

\item $\omega=18{,}8496\,\mathrm{s^{-1}}$.

\item $\delta=0{,}00103006\,\mathrm{m}$.

\item $|u|=0{,}0689916\,\mathrm{m/s}$.

\item $\varphi=0{,}77665$.

\item $|\tau_w|=1{,}85346\,\mathrm{Pa}$.

\item $\overline P/A=0{,}0982948\,\mathrm{W/m^2}$.

\end{itemize}

\textbf{Interpretación.} La perturbación penetra menos al aumentar la frecuencia. La fase indicada es un retraso en radianes, no un tiempo.

\textbf{Comprobación.} Sustituir el perfil en $\partial_tu=\nu\partial_{yy}u$ y comprobar $u(0,t)=U\cos\omega t$. La potencia media es positiva.

\textbf{Límites.} No incluye el transitorio de arranque ni una segunda pared a distancia comparable con la penetración.

\end{studiumworked}

\Needspace{.30\textheight}\section{MF-S2-4. Codo con salida a la atmósfera}

\begin{studiumworked}{MF-S2-4 · Intermedio}

\textbf{Enunciado.} Un codo desvía un chorro un ángulo $\theta$ desde $+x$ hacia $+y$, manteniendo la sección. Calcule la fuerza que el líquido ejerce sobre el codo, indicando componentes y signo.

\begin{longtable}{@{}p{.15\linewidth}p{.23\linewidth}p{.54\linewidth}@{}}
\toprule Símbolo & Valor SI & Significado\\\midrule\endhead
$Q$ & $0{,}015\,\mathrm{m^3/s}$ & caudal\\
$D$ & $0{,}035\,\mathrm{m}$ & diámetro\\
$\rho$ & $1000\,\mathrm{kg/m^3}$ & densidad\\
$p$ & $100000\,\mathrm{Pa}$ & presión manométrica de entrada\\
$\theta$ & $1{,}5708$ & ángulo de salida en radianes\\
\bottomrule\end{longtable}

\textbf{Incógnitas.} Se calculan $A,\ V,\ \dot m,\ F_x,\ F_y,\ |\boldsymbol F|$ y se interpreta el resultado bajo las hipótesis indicadas.

\textbf{Hipótesis.} Balance estacionario con perfiles uniformes, salida atmosférica y peso del volumen de líquido omitido. Las presiones se expresan respecto de la atmósfera.

\textbf{Modelo y ecuaciones.}

\[\fitmath{\boldsymbol R+(pA,0)=\dot m V(\cos\theta-1,\sin\theta),\quad \boldsymbol F_{liquido\to codo}=-\boldsymbol R.}\]

\textbf{Desarrollo.}\begin{enumerate}

\item Escoger el volumen de control interior al codo y calcular velocidad y caudal másico.

\item La presión de entrada empuja hacia $+x$; no se suma una fuerza de presión manométrica en la salida.

\item Obtener primero pared sobre líquido y después invertir el signo para contestar la pregunta.

\end{enumerate}

\textbf{Cálculo de cada magnitud.}

\begin{enumerate}

\item Obtener $A$ con la relación del modelo y los valores anteriores.

\[\fitmath{A=\frac{\pi\cdot {D}^{2}}{4}}\]

\[\fitmath{A=\frac{\pi\cdot {(0{,}035)}^{2}}{4}\approx 0{,}0009621128\,\mathrm{m^{2}}}\]

\item Obtener $V$ con la relación del modelo y los valores anteriores.

\[\fitmath{V=\frac{Q}{A}}\]

\[\fitmath{V=\frac{(0{,}015)}{(0{,}0009621127502)}\approx 15{,}59069\,\mathrm{m/s}}\]

\item Obtener $\dot m$ con la relación del modelo y los valores anteriores.

\[\fitmath{\mathrm{md}=\rho\cdot Q}\]

\[\fitmath{\mathrm{md}=(1000)\cdot (0{,}015)\approx 15\,\mathrm{kg/s}}\]

\item Obtener $F_x$ con la relación del modelo y los valores anteriores.

\[\fitmath{\mathrm{Fx}=p\cdot A-\mathrm{md}\cdot V\cdot \left(\cos\left(\theta\right)-1\right)}\]

\[\fitmath{\mathrm{Fx}=(100000)\cdot (0{,}0009621127502)-(15)\cdot (15{,}5906883)\cdot \left(\cos\left((1{,}570796327)\right)-1\right)\approx 330{,}0716\,\mathrm{N}}\]

\item Obtener $F_y$ con la relación del modelo y los valores anteriores.

\[\fitmath{\mathrm{Fy}=-\mathrm{md}\cdot V\cdot \sin\left(\theta\right)}\]

\[\fitmath{\mathrm{Fy}=-(15)\cdot (15{,}5906883)\cdot \sin\left((1{,}570796327)\right)\approx -233{,}8603\,\mathrm{N}}\]

\item Obtener $|\boldsymbol F|$ con la relación del modelo y los valores anteriores.

\[\fitmath{\mathrm{Fmag}=\sqrt{{\mathrm{Fx}}^{2}+{\mathrm{Fy}}^{2}}}\]

\[\fitmath{\mathrm{Fmag}=\sqrt{{(330{,}0715996)}^{2}+{(-233{,}8603245)}^{2}}\approx 404{,}5218\,\mathrm{N}}\]

\end{enumerate}

\textbf{Resultados.}\begin{itemize}

\item $A=0{,}000962113\,\mathrm{m^2}$.

\item $V=15{,}5907\,\mathrm{m/s}$.

\item $\dot m=15\,\mathrm{kg/s}$.

\item $F_x=330{,}072\,\mathrm{N}$.

\item $F_y=-233{,}86\,\mathrm{N}$.

\item $|\boldsymbol F|=404{,}522\,\mathrm{N}$.

\end{itemize}

\textbf{Interpretación.} La reacción sobre el soporte no se obtiene sólo multiplicando caudal por cambio de velocidad: la presión de entrada también interviene.

\textbf{Comprobación.} Con $\theta=0$ el cambio de momento desaparece y la fuerza axial queda $pA$. Comprobar cada componente por separado.

\textbf{Límites.} Para perfiles no uniformes se necesitan coeficientes de cantidad de movimiento; el peso o las presiones exteriores se incorporan según la frontera elegida.

\end{studiumworked}

\Needspace{.30\textheight}\section{MF-S2-5. Darcy, Fanning y Poiseuille}

\begin{studiumworked}{MF-S2-5 · Intermedio}

\textbf{Enunciado.} Calcule pérdidas en un tubo con la ley laminar, indicando los factores de Darcy y Fanning. Verifique el resultado mediante Poiseuille y explique el error de mezclar las dos convenciones.

\begin{longtable}{@{}p{.15\linewidth}p{.23\linewidth}p{.54\linewidth}@{}}
\toprule Símbolo & Valor SI & Significado\\\midrule\endhead
$\rho$ & $900\,\mathrm{kg/m^3}$ & densidad\\
$\mu$ & $0{,}009\,\mathrm{Pa\,s}$ & viscosidad\\
$D$ & $0{,}012\,\mathrm{m}$ & diámetro\\
$V$ & $0{,}8\,\mathrm{m/s}$ & velocidad media\\
$L$ & $3\,\mathrm{m}$ & longitud\\
\bottomrule\end{longtable}

\textbf{Incógnitas.} Se calculan $\mathrm{Re}_D,\ f_D,\ f_F,\ \Delta p,\ \Delta p_{Pois}$ y se interpreta el resultado bajo las hipótesis indicadas.

\textbf{Hipótesis.} Tubo circular horizontal, flujo laminar newtoniano plenamente desarrollado. Se omiten pérdidas locales y de entrada.

\textbf{Modelo y ecuaciones.}

\[\fitmath{f_D=64/\mathrm{Re},\quad f_F=f_D/4,\quad \Delta p=f_D(L/D)\rho V^2/2=32\mu LV/D^2.}\]

\textbf{Desarrollo.}\begin{enumerate}

\item Calcular Reynolds con velocidad media y diámetro.

\item Aplicar el factor de Darcy en su ecuación de caída de presión.

\item Obtener la misma caída a partir del balance laminar; introducir $f_F$ sin su factor cuatro daría una pérdida cuatro veces menor.

\end{enumerate}

\textbf{Cálculo de cada magnitud.}

\begin{enumerate}

\item Obtener $\mathrm{Re}_D$ con la relación del modelo y los valores anteriores.

\[\fitmath{\mathrm{Re}=\frac{\rho\cdot V\cdot D}{\mu}}\]

\[\fitmath{\mathrm{Re}=\frac{(900)\cdot (0{,}8)\cdot (0{,}012)}{(0{,}009)}\approx 960\,}\]

\item Obtener $f_D$ con la relación del modelo y los valores anteriores.

\[\fitmath{\mathrm{fD}=\frac{64}{\mathrm{Re}}}\]

\[\fitmath{\mathrm{fD}=\frac{64}{(960)}\approx 0{,}06666667\,}\]

\item Obtener $f_F$ con la relación del modelo y los valores anteriores.

\[\fitmath{\mathrm{fF}=\frac{16}{\mathrm{Re}}}\]

\[\fitmath{\mathrm{fF}=\frac{16}{(960)}\approx 0{,}01666667\,}\]

\item Obtener $\Delta p$ con la relación del modelo y los valores anteriores.

\[\fitmath{\mathrm{dp}=\frac{\mathrm{fD}\cdot \frac{L}{D}\cdot \rho\cdot {V}^{2}}{2}}\]

\[\fitmath{\mathrm{dp}=\frac{(0{,}06666666667)\cdot \frac{(3)}{(0{,}012)}\cdot (900)\cdot {(0{,}8)}^{2}}{2}\approx 4800\,\mathrm{Pa}}\]

\item Obtener $\Delta p_{Pois}$ con la relación del modelo y los valores anteriores.

\[\fitmath{\mathrm{dp}_{2}=\frac{32\cdot \mu\cdot L\cdot V}{{D}^{2}}}\]

\[\fitmath{\mathrm{dp}_{2}=\frac{32\cdot (0{,}009)\cdot (3)\cdot (0{,}8)}{{(0{,}012)}^{2}}\approx 4800\,\mathrm{Pa}}\]

\end{enumerate}

\textbf{Resultados.}\begin{itemize}

\item $\mathrm{Re}_D=960$.

\item $f_D=0{,}0666667$.

\item $f_F=0{,}0166667$.

\item $\Delta p=4800\,\mathrm{Pa}$.

\item $\Delta p_{Pois}=4800\,\mathrm{Pa}$.

\end{itemize}

\textbf{Interpretación.} La notación del factor de fricción forma parte de la definición del modelo; una tabla sin convenio puede inducir errores grandes.

\textbf{Comprobación.} Comprobar $f_D=4f_F$ y la igualdad de las dos rutas de presión.

\textbf{Límites.} La transición no se toma como una frontera universal exacta; la ley $64/\mathrm{Re}$ no se extiende a régimen turbulento.

\end{studiumworked}

\studiumchaptercheck

\chapter{Soluciones del simulacro 3}

\objetivos{Comparar el planteamiento antes de consultar los resultados. La rúbrica valora el modelo, no sólo la cifra final.}

\Needspace{.30\textheight}\section{MF-S3-1. Viscosidad a partir de un ensayo de placas}

\begin{studiumworked}{MF-S3-1 · Básico}

\textbf{Enunciado.} Una placa arrastra una capa de líquido. A partir de la fuerza tangencial, obtenga viscosidad dinámica y cinemática, esfuerzo y potencia disipada. Calcule un Reynolds basado en el espesor.

\begin{longtable}{@{}p{.15\linewidth}p{.23\linewidth}p{.54\linewidth}@{}}
\toprule Símbolo & Valor SI & Significado\\\midrule\endhead
$F$ & $4{,}2\,\mathrm{N}$ & fuerza tangencial\\
$A$ & $0{,}14\,\mathrm{m^2}$ & área mojada\\
$h$ & $0{,}0015\,\mathrm{m}$ & separación\\
$U$ & $0{,}5\,\mathrm{m/s}$ & velocidad de la placa\\
$\rho$ & $950\,\mathrm{kg/m^3}$ & densidad\\
\bottomrule\end{longtable}

\textbf{Incógnitas.} Se calculan $\tau,\ \mu,\ \nu,\ P,\ \mathrm{Re}_h$ y se interpreta el resultado bajo las hipótesis indicadas.

\textbf{Hipótesis.} Fluido newtoniano, perfil lineal, no deslizamiento, placa opuesta fija y gradiente de presión nulo. Fuerza neta de una sola cara mojada.

\textbf{Modelo y ecuaciones.}

\[\fitmath{\tau=F/A=\mu U/h,\quad \mu=Fh/(AU),\quad \nu=\mu/\rho,\quad P=FU.}\]

\textbf{Desarrollo.}\begin{enumerate}

\item Convertir fuerza en esfuerzo usando el área mojada correcta.

\item Aplicar la ley constitutiva al gradiente uniforme $U/h$.

\item La potencia de la placa se disipa por viscosidad: por unidad de volumen es $\mu(U/h)^2$.

\end{enumerate}

\textbf{Cálculo de cada magnitud.}

\begin{enumerate}

\item Obtener $\tau$ con la relación del modelo y los valores anteriores.

\[\fitmath{\tau=\frac{F}{A}}\]

\[\fitmath{\tau=\frac{(4{,}2)}{(0{,}14)}\approx 30\,\mathrm{Pa}}\]

\item Obtener $\mu$ con la relación del modelo y los valores anteriores.

\[\fitmath{\mu=\frac{F\cdot h}{A\cdot U}}\]

\[\fitmath{\mu=\frac{(4{,}2)\cdot (0{,}0015)}{(0{,}14)\cdot (0{,}5)}\approx 0{,}09\,\mathrm{Pa\,s}}\]

\item Obtener $\nu$ con la relación del modelo y los valores anteriores.

\[\fitmath{\nu=\frac{\mu}{\rho}}\]

\[\fitmath{\nu=\frac{(0{,}09)}{(950)}\approx 9{,}473684\times10^{-5}\,\mathrm{m^{2}/s}}\]

\item Obtener $P$ con la relación del modelo y los valores anteriores.

\[\fitmath{\mathrm{power}=F\cdot U}\]

\[\fitmath{\mathrm{power}=(4{,}2)\cdot (0{,}5)\approx 2{,}1\,\mathrm{W}}\]

\item Obtener $\mathrm{Re}_h$ con la relación del modelo y los valores anteriores.

\[\fitmath{\mathrm{Re}=\frac{\rho\cdot U\cdot h}{\mu}}\]

\[\fitmath{\mathrm{Re}=\frac{(950)\cdot (0{,}5)\cdot (0{,}0015)}{(0{,}09)}\approx 7{,}916667\,}\]

\end{enumerate}

\textbf{Resultados.}\begin{itemize}

\item $\tau=30\,\mathrm{Pa}$.

\item $\mu=0{,}09\,\mathrm{Pa\,s}$.

\item $\nu=9{,}47368\times10^{-5}\,\mathrm{m^2/s}$.

\item $P=2{,}1\,\mathrm{W}$.

\item $\mathrm{Re}_h=7{,}91667$.

\end{itemize}

\textbf{Interpretación.} Viscosidad dinámica y cinemática tienen dimensiones distintas. Un Reynolds pequeño es coherente con este caso, pero no sustituye la comprobación de las condiciones de contorno.

\textbf{Comprobación.} Integrar $\mu(U/h)^2$ sobre el volumen $Ah$ y recuperar $FU$.

\textbf{Límites.} Se omiten efectos de borde, calentamiento y dependencia de la viscosidad con temperatura o cizalla.

\end{studiumworked}

\Needspace{.30\textheight}\section{MF-S3-2. Máxima eficiencia con una polar parabólica}

\begin{studiumworked}{MF-S3-2 · Intermedio}

\textbf{Enunciado.} Un avión en vuelo horizontal tiene $C_D=C_{D0}+KC_L^2$. Halle el $C_L$, velocidad y resistencia de máximo $L/D$. Separe la contribución inducida $K_i=1/(\pi eAR)$ de otras contribuciones cuadráticas.

\begin{longtable}{@{}p{.15\linewidth}p{.23\linewidth}p{.54\linewidth}@{}}
\toprule Símbolo & Valor SI & Significado\\\midrule\endhead
$C_{D0}$ & $0{,}028$ & coeficiente constante\\
$K$ & $0{,}052$ & coeficiente cuadrático total\\
$e$ & $0{,}9$ & factor inducido dado\\
$AR$ & $9$ & alargamiento\\
$W$ & $5000\,\mathrm{N}$ & peso\\
$\rho$ & $1{,}1\,\mathrm{kg/m^3}$ & densidad\\
$S$ & $14\,\mathrm{m^2}$ & superficie\\
\bottomrule\end{longtable}

\textbf{Incógnitas.} Se calculan $C_{L,*},\ V_*,\ D_{min},\ K_i,\ K-K_i,\ (L/D)_{max}$ y se interpreta el resultado bajo las hipótesis indicadas.

\textbf{Hipótesis.} Vuelo estacionario horizontal, $L=W$, polar válida en el intervalo de sustentación empleado; datos sintéticos sin atribución a un avión.

\textbf{Modelo y ecuaciones.}

\[\fitmath{D/W=C_{D0}/C_L+KC_L,\quad C_{L,*}=\sqrt{C_{D0}/K},\quad V_*=\sqrt{2W/(\rho SC_{L,*})},\quad D_{min}=2W\sqrt{C_{D0}K}.}\]

\textbf{Desarrollo.}\begin{enumerate}

\item Derivar la resistencia respecto de $C_L>0$ con peso fijo y comprobar mínimo.

\item Obtener la velocidad a partir del equilibrio de sustentación.

\item Calcular $K_i$ y el resto $K-K_i$. La igualdad óptima es entre término constante y término cuadrático total, no necesariamente resistencia inducida sola.

\end{enumerate}

\textbf{Cálculo de cada magnitud.}

\begin{enumerate}

\item Obtener $C_{L,*}$ con la relación del modelo y los valores anteriores.

\[\fitmath{C_L=\sqrt{\frac{\mathrm{cd}_{0}}{K}}}\]

\[\fitmath{C_L=\sqrt{\frac{(0{,}028)}{(0{,}052)}}\approx 0{,}7337994\,}\]

\item Obtener $V_*$ con la relación del modelo y los valores anteriores.

\[\fitmath{V=\sqrt{\frac{2\cdot W}{\rho\cdot S\cdot C_L}}}\]

\[\fitmath{V=\sqrt{\frac{2\cdot (5000)}{(1{,}1)\cdot (14)\cdot (0{,}7337993857)}}\approx 29{,}74753\,\mathrm{m/s}}\]

\item Obtener $D_{min}$ con la relación del modelo y los valores anteriores.

\[\fitmath{\mathrm{Dmin}=2\cdot W\cdot \sqrt{\mathrm{cd}_{0}\cdot K}}\]

\[\fitmath{\mathrm{Dmin}=2\cdot (5000)\cdot \sqrt{(0{,}028)\cdot (0{,}052)}\approx 381{,}5757\,\mathrm{N}}\]

\item Obtener $K_i$ con la relación del modelo y los valores anteriores.

\[\fitmath{\mathrm{Ki}=\frac{1}{\pi\cdot e\cdot \mathrm{AR}}}\]

\[\fitmath{\mathrm{Ki}=\frac{1}{\pi\cdot (0{,}9)\cdot (9)}\approx 0{,}03929752\,}\]

\item Obtener $K-K_i$ con la relación del modelo y los valores anteriores.

\[\fitmath{\mathrm{rest}=K-\mathrm{Ki}}\]

\[\fitmath{\mathrm{rest}=(0{,}052)-(0{,}03929751681)\approx 0{,}01270248\,}\]

\item Obtener $(L/D)_{max}$ con la relación del modelo y los valores anteriores.

\[\fitmath{\mathrm{LD}=\frac{W}{\mathrm{Dmin}}}\]

\[\fitmath{\mathrm{LD}=\frac{(5000)}{(381{,}5756806)}\approx 13{,}10356\,}\]

\end{enumerate}

\textbf{Resultados.}\begin{itemize}

\item $C_{L,*}=0{,}733799$.

\item $V_*=29{,}7475\,\mathrm{m/s}$.

\item $D_{min}=381{,}576\,\mathrm{N}$.

\item $K_i=0{,}0392975$.

\item $K-K_i=0{,}0127025$.

\item $(L/D)_{max}=13{,}1036$.

\end{itemize}

\textbf{Interpretación.} Atribuir todo el coeficiente cuadrático a inducción sólo es válido cuando el modelo de la polar lo establece.

\textbf{Comprobación.} La segunda derivada $2C_{D0}/C_L^3$ es positiva; verificar $KC_{L,*}^2=C_{D0}$.

\textbf{Límites.} No se han comprobado pérdida, límites estructurales, variación de polar con Reynolds o compresibilidad.

\end{studiumworked}

\Needspace{.30\textheight}\section{MF-S3-3. Fuente y vórtice en un anillo}

\begin{studiumworked}{MF-S3-3 · Avanzado}

\textbf{Enunciado.} En un dominio anular se prescribe $u_r=C/r$, $u_\theta=K/r$. Halle caudal radial por longitud, circulación y aumento de presión entre dos radios. Explique por qué puede existir circulación con rotacional local nulo.

\begin{longtable}{@{}p{.15\linewidth}p{.23\linewidth}p{.54\linewidth}@{}}
\toprule Símbolo & Valor SI & Significado\\\midrule\endhead
$C$ & $0{,}004\,\mathrm{m^2/s}$ & intensidad radial\\
$K$ & $0{,}008\,\mathrm{m^2/s}$ & intensidad azimutal\\
$r_1$ & $0{,}12\,\mathrm{m}$ & radio interior\\
$r_2$ & $0{,}24\,\mathrm{m}$ & radio exterior\\
$\rho$ & $1000\,\mathrm{kg/m^3}$ & densidad\\
\bottomrule\end{longtable}

\textbf{Incógnitas.} Se calculan $Q^{\prime},\ \Gamma,\ p_2-p_1$ y se interpreta el resultado bajo las hipótesis indicadas.

\textbf{Hipótesis.} Flujo ideal estacionario bidimensional en $r_1\le r\le r_2$, sin incluir la singularidad del origen; densidad constante.

\textbf{Modelo y ecuaciones.}

\[\fitmath{Q^{\prime}=2\pi C,\quad \Gamma=2\pi K,\quad \nabla\cdot\boldsymbol u=0,\quad \nabla\times\boldsymbol u=0\ (r>0),\quad p_2-p_1=\tfrac12\rho(C^2+K^2)(r_1^{-2}-r_2^{-2}).}\]

\textbf{Desarrollo.}\begin{enumerate}

\item Integrar las componentes normal y tangencial sobre un círculo.

\item Calcular operadores en coordenadas cilíndricas y conservar los factores $r$.

\item Aplicar Bernoulli en el campo irrotacional. Un contorno que rodea el agujero no limita una superficie enteramente contenida en el dominio del campo.

\end{enumerate}

\textbf{Cálculo de cada magnitud.}

\begin{enumerate}

\item Obtener $Q^{\prime}$ con la relación del modelo y los valores anteriores.

\[\fitmath{\mathrm{flux}=2\cdot \pi\cdot C}\]

\[\fitmath{\mathrm{flux}=2\cdot \pi\cdot (0{,}004)\approx 0{,}02513274\,\mathrm{m^{2}/s}}\]

\item Obtener $\Gamma$ con la relación del modelo y los valores anteriores.

\[\fitmath{\mathrm{circ}=2\cdot \pi\cdot K}\]

\[\fitmath{\mathrm{circ}=2\cdot \pi\cdot (0{,}008)\approx 0{,}05026548\,\mathrm{m^{2}/s}}\]

\item Obtener $p_2-p_1$ con la relación del modelo y los valores anteriores.

\[\fitmath{\mathrm{dp}=\frac{\rho\cdot \left({C}^{2}+{K}^{2}\right)\cdot \left(\frac{1}{{r_{1}}^{2}}-\frac{1}{{r_{2}}^{2}}\right)}{2}}\]

\[\fitmath{\mathrm{dp}=\frac{(1000)\cdot \left({(0{,}004)}^{2}+{(0{,}008)}^{2}\right)\cdot \left(\frac{1}{{(0{,}12)}^{2}}-\frac{1}{{(0{,}24)}^{2}}\right)}{2}\approx 2{,}083333\,\mathrm{Pa}}\]

\end{enumerate}

\textbf{Resultados.}\begin{itemize}

\item $Q^{\prime}=0{,}0251327\,\mathrm{m^2/s}$.

\item $\Gamma=0{,}0502655\,\mathrm{m^2/s}$.

\item $p_2-p_1=2{,}08333\,\mathrm{Pa}$.

\end{itemize}

\textbf{Interpretación.} La topología del dominio es parte de las hipótesis del teorema: no se puede aplicar Stokes atravesando la singularidad excluida.

\textbf{Comprobación.} Derivar la presión: $dp/dr=\rho(C^2+K^2)/r^3$. Para $K=0$ se conserva el caudal radial y desaparece circulación.

\textbf{Límites.} No satisface no deslizamiento en una pared sólida anular y no modela el núcleo de un vórtice real.

\end{studiumworked}

\Needspace{.30\textheight}\section{MF-S3-4. Vaciado cuasiestacionario de un depósito}

\begin{studiumworked}{MF-S3-4 · Intermedio}

\textbf{Enunciado.} Un depósito de sección constante descarga por un orificio. Deduzca el tiempo para bajar entre los niveles dados y el caudal inicial, usando el coeficiente de descarga proporcionado.

\begin{longtable}{@{}p{.15\linewidth}p{.23\linewidth}p{.54\linewidth}@{}}
\toprule Símbolo & Valor SI & Significado\\\midrule\endhead
$A_t$ & $1\,\mathrm{m^2}$ & sección del depósito\\
$A_o$ & $0{,}0005\,\mathrm{m^2}$ & área del orificio\\
$C_d$ & $0{,}62$ & coeficiente de descarga dado\\
$h_0$ & $2{,}5\,\mathrm{m}$ & nivel inicial sobre orificio\\
$h_1$ & $0{,}8\,\mathrm{m}$ & nivel final sobre orificio\\
\bottomrule\end{longtable}

\textbf{Incógnitas.} Se calculan $t,\ Q_0,\ A_o/A_t$ y se interpreta el resultado bajo las hipótesis indicadas.

\textbf{Hipótesis.} Superficie y salida atmosféricas, líquido incompresible, orificio pequeño respecto del depósito y coeficiente constante entre los niveles dados.

\textbf{Modelo y ecuaciones.}

\[\fitmath{A_t\dot h=-C_dA_o\sqrt{2gh},\quad t=\frac{2A_t}{C_dA_o\sqrt{2g}}(\sqrt{h_0}-\sqrt{h_1}).}\]

\textbf{Desarrollo.}\begin{enumerate}

\item Conectar Bernoulli ideal con la corrección empírica $C_d$ suministrada; no deducirla por análisis dimensional.

\item Usar continuidad para relacionar caudal con descenso de nivel.

\item Separar variables e integrar con los límites en el orden que haga positivo el tiempo.

\end{enumerate}

\textbf{Cálculo de cada magnitud.}

\begin{enumerate}

\item Obtener $t$ con la relación del modelo y los valores anteriores.

\[\fitmath{\mathrm{time}=\frac{2\cdot \mathrm{At}}{\mathrm{Cd}\cdot \mathrm{Ao}\cdot \sqrt{2\cdot g}}\cdot \left(\sqrt{h_{0}}-\sqrt{h_{1}}\right)}\]

\[\fitmath{\mathrm{time}=\frac{2\cdot (1)}{(0{,}62)\cdot (0{,}0005)\cdot \sqrt{2\cdot (9{,}81)}}\cdot \left(\sqrt{(2{,}5)}-\sqrt{(0{,}8)}\right)\approx 1000{,}215\,\mathrm{s}}\]

\item Obtener $Q_0$ con la relación del modelo y los valores anteriores.

\[\fitmath{Q_{0}=\mathrm{Cd}\cdot \mathrm{Ao}\cdot \sqrt{2\cdot g\cdot h_{0}}}\]

\[\fitmath{Q_{0}=(0{,}62)\cdot (0{,}0005)\cdot \sqrt{2\cdot (9{,}81)\cdot (2{,}5)}\approx 0{,}002171107\,\mathrm{m^{3}/s}}\]

\item Obtener $A_o/A_t$ con la relación del modelo y los valores anteriores.

\[\fitmath{\mathrm{ratio}=\frac{\mathrm{Ao}}{\mathrm{At}}}\]

\[\fitmath{\mathrm{ratio}=\frac{(0{,}0005)}{(1)}\approx 0{,}0005\,}\]

\end{enumerate}

\textbf{Resultados.}\begin{itemize}

\item $t=1000{,}21\,\mathrm{s}$.

\item $Q_0=0{,}00217111\,\mathrm{m^3/s}$.

\item $A_o/A_t=0{,}0005$.

\end{itemize}

\textbf{Interpretación.} El caudal decrece al bajar el nivel. Sustituirlo por el inicial durante todo el proceso subestima el tiempo.

\textbf{Comprobación.} Integrar la ecuación diferencial numéricamente y localizar $h=h_1$. Para $h_1=h_0$ el tiempo es cero.

\textbf{Límites.} No extrapolar al vaciado final si cambian la geometría efectiva, la ventilación o el coeficiente de descarga.

\end{studiumworked}

\Needspace{.30\textheight}\section{MF-S3-5. Tracción y disipación en un campo lineal}

\begin{studiumworked}{MF-S3-5 · Intermedio}

\textbf{Enunciado.} Para $u=ax+by$, $v=-ay$, $w=0$, calcule tensor de esfuerzos, tracción normal y tangencial sobre el plano de normal $(1,1,0)/\sqrt2$ y disipación viscosa. Use tangente $(-1,1,0)/\sqrt2$.

\begin{longtable}{@{}p{.15\linewidth}p{.23\linewidth}p{.54\linewidth}@{}}
\toprule Símbolo & Valor SI & Significado\\\midrule\endhead
$a$ & $1{,}5\,\mathrm{s^{-1}}$ & extensión\\
$b$ & $4\,\mathrm{s^{-1}}$ & cizalla\\
$\mu$ & $0{,}3\,\mathrm{Pa\,s}$ & viscosidad\\
$p$ & $1500\,\mathrm{Pa}$ & presión\\
\bottomrule\end{longtable}

\textbf{Incógnitas.} Se calculan $\sigma_{nn},\ \tau_{tn},\ \Phi$ y se interpreta el resultado bajo las hipótesis indicadas.

\textbf{Hipótesis.} Fluido newtoniano incompresible isotrópico, esfuerzo $\boldsymbol\sigma=-p\boldsymbol I+2\mu\boldsymbol D$. La tracción se define como $\boldsymbol\sigma\boldsymbol n$.

\textbf{Modelo y ecuaciones.}

\[\fitmath{D=\begin{pmatrix}a&b/2&0\\b/2&-a&0\\0&0&0\end{pmatrix},\quad \sigma_{nn}=-p+\mu b,\quad\tau_{tn}=-2\mu a,\quad\Phi=2\mu D:D=\mu(4a^2+b^2).}\]

\textbf{Desarrollo.}\begin{enumerate}

\item Simetrizar el gradiente y comprobar traza nula.

\item Multiplicar tensor por normal y proyectar el resultado en la normal y la tangente orientada dadas.

\item Contraer el tensor de deformación consigo mismo, contando los dos términos fuera de la diagonal.

\end{enumerate}

\textbf{Cálculo de cada magnitud.}

\begin{enumerate}

\item Obtener $\sigma_{nn}$ con la relación del modelo y los valores anteriores.

\[\fitmath{\mathrm{sn}=-p+\mu\cdot b}\]

\[\fitmath{\mathrm{sn}=-(1500)+(0{,}3)\cdot (4)\approx -1498{,}8\,\mathrm{Pa}}\]

\item Obtener $\tau_{tn}$ con la relación del modelo y los valores anteriores.

\[\fitmath{\mathrm{st}=-2\cdot \mu\cdot a}\]

\[\fitmath{\mathrm{st}=-2\cdot (0{,}3)\cdot (1{,}5)\approx -0{,}9\,\mathrm{Pa}}\]

\item Obtener $\Phi$ con la relación del modelo y los valores anteriores.

\[\fitmath{\varphi=\mu\cdot \left(4\cdot {a}^{2}+{b}^{2}\right)}\]

\[\fitmath{\varphi=(0{,}3)\cdot \left(4\cdot {(1{,}5)}^{2}+{(4)}^{2}\right)\approx 7{,}5\,\mathrm{W/m^{3}}}\]

\end{enumerate}

\textbf{Resultados.}\begin{itemize}

\item $\sigma_{nn}=-1498{,}8\,\mathrm{Pa}$.

\item $\tau_{tn}=-0{,}9\,\mathrm{Pa}$.

\item $\Phi=7{,}5\,\mathrm{W/m^3}$.

\end{itemize}

\textbf{Interpretación.} La presión contribuye a la tracción normal, mientras la parte viscosa puede tener componentes normales y tangenciales.

\textbf{Comprobación.} La disipación debe ser no negativa para $\mu>0$ y desaparecer si $D=0$. Invertir la tangente cambia el signo de $\tau_{tn}$.

\textbf{Límites.} No se usa esta ley constitutiva para fluidos no newtonianos ni para compresión volumétrica arbitraria.

\end{studiumworked}

\studiumchaptercheck

\chapter{Soluciones del simulacro 4}

\objetivos{Comparar el planteamiento antes de consultar los resultados. La rúbrica valora el modelo, no sólo la cifra final.}

\Needspace{.30\textheight}\section{MF-S4-1. Oscilador amortiguado desde Lagrange}

\begin{studiumworked}{MF-S4-1 · Intermedio}

\textbf{Enunciado.} Una masa unida a un muelle y un amortiguador parte de una elongación sin velocidad. Obtenga su ecuación, determine el régimen y calcule frecuencia amortiguada y amplitud de la envolvente al tiempo dado.

\begin{longtable}{@{}p{.15\linewidth}p{.23\linewidth}p{.54\linewidth}@{}}
\toprule Símbolo & Valor SI & Significado\\\midrule\endhead
$m$ & $3\,\mathrm{kg}$ & masa\\
$k$ & $120\,\mathrm{N/m}$ & rigidez\\
$b$ & $6\,\mathrm{kg/s}$ & amortiguamiento\\
$x_0$ & $0{,}04\,\mathrm{m}$ & elongación inicial\\
$t$ & $1\,\mathrm{s}$ & tiempo\\
\bottomrule\end{longtable}

\textbf{Incógnitas.} Se calculan $\omega_n,\ \zeta,\ \omega_d,\ A_{env},\ x$ y se interpreta el resultado bajo las hipótesis indicadas.

\textbf{Hipótesis.} Un grado de libertad, muelle lineal y fuerza disipativa $-b\dot x$; no hay excitación exterior.

\textbf{Modelo y ecuaciones.}

\[\fitmath{L=\tfrac12m\dot x^2-\tfrac12kx^2,\quad m\ddot x+b\dot x+kx=0,\quad \zeta=\frac{b}{2\sqrt{km}},\quad \omega_d=\sqrt{k/m-b^2/(4m^2)}.}\]

\textbf{Desarrollo.}\begin{enumerate}

\item Aplicar Euler--Lagrange con fuerza generalizada $-b\dot x$.

\item Comparar $\zeta$ con uno y resolver las raíces de la ecuación característica.

\item La solución con velocidad inicial nula es $x=x_0e^{-bt/(2m)}[\cos(\omega_dt)+b\sin(\omega_dt)/(2m\omega_d)]$. La envolvente exacta de esta expresión incluye el factor de fase.

\end{enumerate}

\textbf{Cálculo de cada magnitud.}

\begin{enumerate}

\item Obtener $\omega_n$ con la relación del modelo y los valores anteriores.

\[\fitmath{\mathrm{wn}=\sqrt{\frac{k}{m}}}\]

\[\fitmath{\mathrm{wn}=\sqrt{\frac{(120)}{(3)}}\approx 6{,}324555\,\mathrm{1/s}}\]

\item Obtener $\zeta$ con la relación del modelo y los valores anteriores.

\[\fitmath{\zeta=\frac{b}{2\cdot \sqrt{k\cdot m}}}\]

\[\fitmath{\zeta=\frac{(6)}{2\cdot \sqrt{(120)\cdot (3)}}\approx 0{,}1581139\,}\]

\item Obtener $\omega_d$ con la relación del modelo y los valores anteriores.

\[\fitmath{\mathrm{wd}=\sqrt{\frac{k}{m}-\frac{{b}^{2}}{4\cdot {m}^{2}}}}\]

\[\fitmath{\mathrm{wd}=\sqrt{\frac{(120)}{(3)}-\frac{{(6)}^{2}}{4\cdot {(3)}^{2}}}\approx 6{,}244998\,\mathrm{1/s}}\]

\item Obtener $A_{env}$ con la relación del modelo y los valores anteriores.

\[\fitmath{\mathrm{env}=x_{0}\cdot \sqrt{1+{\left(\frac{b}{2\cdot m\cdot \mathrm{wd}}\right)}^{2}}\cdot \exp\left(\frac{-b\cdot t}{2\cdot m}\right)}\]

\[\fitmath{\mathrm{env}=(0{,}04)\cdot \sqrt{1+{\left(\frac{(6)}{2\cdot (3)\cdot (6{,}244997998)}\right)}^{2}}\cdot \exp\left(\frac{-(6)\cdot (1)}{2\cdot (3)}\right)\approx 0{,}01490264\,\mathrm{m}}\]

\item Obtener $x$ con la relación del modelo y los valores anteriores.

\[\fitmath{x=x_{0}\cdot \exp\left(\frac{-b\cdot t}{2\cdot m}\right)\cdot \left(\cos\left(\mathrm{wd}\cdot t\right)+\frac{b\cdot \sin\left(\mathrm{wd}\cdot t\right)}{2\cdot m\cdot \mathrm{wd}}\right)}\]

\[\fitmath{x=(0{,}04)\cdot \exp\left(\frac{-(6)\cdot (1)}{2\cdot (3)}\right)\cdot \left(\cos\left((6{,}244997998)\cdot (1)\right)+\frac{(6)\cdot \sin\left((6{,}244997998)\cdot (1)\right)}{2\cdot (3)\cdot (6{,}244997998)}\right)\approx 0{,}01461449\,\mathrm{m}}\]

\end{enumerate}

\textbf{Resultados.}\begin{itemize}

\item $\omega_n=6{,}32456\,\mathrm{s^{-1}}$.

\item $\zeta=0{,}158114$.

\item $\omega_d=6{,}245\,\mathrm{s^{-1}}$.

\item $A_{env}=0{,}0149026\,\mathrm{m}$.

\item $x=0{,}0146145\,\mathrm{m}$.

\end{itemize}

\textbf{Interpretación.} La envolvente decrece exponencialmente, pero no es idéntica a la posición instantánea ni, con estas condiciones iniciales, a $x_0e^{-bt/(2m)}$.

\textbf{Comprobación.} $\dot E=-b\dot x^2\le0$; al hacer $b\to0$ se recupera la oscilación conservativa.

\textbf{Límites.} La linealidad del muelle y del amortiguador limita la extrapolación a amplitudes grandes.

\end{studiumworked}

\Needspace{.30\textheight}\section{MF-S4-2. Reynolds pequeño con inercia no estacionaria}

\begin{studiumworked}{MF-S4-2 · Avanzado}

\textbf{Enunciado.} Un flujo de baja velocidad está forzado en un tiempo muy corto. Calcule Reynolds, Strouhal y su producto para decidir qué términos inerciales pueden despreciarse. Estime la escala viscosa de presión.

\begin{longtable}{@{}p{.15\linewidth}p{.23\linewidth}p{.54\linewidth}@{}}
\toprule Símbolo & Valor SI & Significado\\\midrule\endhead
$\rho$ & $1000\,\mathrm{kg/m^3}$ & densidad\\
$\mu$ & $1\,\mathrm{Pa\,s}$ & viscosidad\\
$L$ & $0{,}0012\,\mathrm{m}$ & longitud\\
$U$ & $0{,}0002\,\mathrm{m/s}$ & velocidad\\
$T$ & $0{,}0002\,\mathrm{s}$ & tiempo impuesto\\
\bottomrule\end{longtable}

\textbf{Incógnitas.} Se calculan $\mathrm{Re},\ \mathrm{St},\ \mathrm{Re}\,\mathrm{St},\ p_\nu$ y se interpreta el resultado bajo las hipótesis indicadas.

\textbf{Hipótesis.} Fluido newtoniano, longitud espacial $L$ y escala temporal externa $T$ independientes; comparación de términos de Navier--Stokes.

\textbf{Modelo y ecuaciones.}

\[\fitmath{\mathrm{Re}=\rho UL/\mu,\quad \mathrm{St}=L/(UT),\quad\mathrm{Re}\,\mathrm{St}=\rho L^2/(\mu T),\quad p_\nu=\mu U/L.}\]

\textbf{Desarrollo.}\begin{enumerate}

\item Comparar inercia convectiva con viscosidad mediante Reynolds.

\item Comparar inercia local con viscosidad: el cociente es Reynolds por Strouhal, no Reynolds solo.

\item La presión viscosa sirve de referencia para escribir la ecuación adimensional; cuando la inercia local domina, puede no ser la presión efectiva dominante.

\end{enumerate}

\textbf{Cálculo de cada magnitud.}

\begin{enumerate}

\item Obtener $\mathrm{Re}$ con la relación del modelo y los valores anteriores.

\[\fitmath{\mathrm{Re}=\frac{\rho\cdot U\cdot L}{\mu}}\]

\[\fitmath{\mathrm{Re}=\frac{(1000)\cdot (0{,}0002)\cdot (0{,}0012)}{(1)}\approx 0{,}00024\,}\]

\item Obtener $\mathrm{St}$ con la relación del modelo y los valores anteriores.

\[\fitmath{\mathrm{St}=\frac{L}{U\cdot T}}\]

\[\fitmath{\mathrm{St}=\frac{(0{,}0012)}{(0{,}0002)\cdot (0{,}0002)}\approx 30000\,}\]

\item Obtener $\mathrm{Re}\,\mathrm{St}$ con la relación del modelo y los valores anteriores.

\[\fitmath{\mathrm{product}=\mathrm{Re}\cdot \mathrm{St}}\]

\[\fitmath{\mathrm{product}=(0{,}00024)\cdot (30000)\approx 7{,}2\,}\]

\item Obtener $p_\nu$ con la relación del modelo y los valores anteriores.

\[\fitmath{\mathrm{pressure}=\frac{\mu\cdot U}{L}}\]

\[\fitmath{\mathrm{pressure}=\frac{(1)\cdot (0{,}0002)}{(0{,}0012)}\approx 0{,}1666667\,\mathrm{Pa}}\]

\end{enumerate}

\textbf{Resultados.}\begin{itemize}

\item $\mathrm{Re}=0{,}00024$.

\item $\mathrm{St}=30000$.

\item $\mathrm{Re}\,\mathrm{St}=7{,}2$.

\item $p_\nu=0{,}166667\,\mathrm{Pa}$.

\end{itemize}

\textbf{Interpretación.} La convección puede ser despreciable mientras la derivada temporal sigue siendo relevante o dominante. No es lícito imponer Stokes estacionario sólo por velocidad pequeña.

\textbf{Comprobación.} Comprobar el producto directamente como $L^2/(\nu T)$. Para tiempos largos ese producto disminuye.

\textbf{Límites.} El análisis ordena términos; no resuelve un campo sin geometría, condiciones de contorno y forma del forzamiento.

\end{studiumworked}

\Needspace{.30\textheight}\section{MF-S4-3. Carro móvil que intercepta un chorro}

\begin{studiumworked}{MF-S4-3 · Intermedio}

\textbf{Enunciado.} Un carro se aleja en la dirección de un chorro y desvía el agua noventa grados en su sistema relativo. Halle caudal capturado, fuerza propulsora y potencia mecánica. Compruebe la energía en el sistema del laboratorio.

\begin{longtable}{@{}p{.15\linewidth}p{.23\linewidth}p{.54\linewidth}@{}}
\toprule Símbolo & Valor SI & Significado\\\midrule\endhead
$V$ & $24\,\mathrm{m/s}$ & velocidad del chorro en laboratorio\\
$U$ & $6\,\mathrm{m/s}$ & velocidad constante del carro\\
$A$ & $0{,}0008\,\mathrm{m^2}$ & sección del chorro\\
$\rho$ & $1000\,\mathrm{kg/m^3}$ & densidad\\
\bottomrule\end{longtable}

\textbf{Incógnitas.} Se calculan $\dot m_c,\ F,\ P,\ \Delta\dot E_c$ y se interpreta el resultado bajo las hipótesis indicadas.

\textbf{Hipótesis.} Chorro libre con $V>U$, sin pérdidas en la desviación relativa; la componente longitudinal relativa de salida es nula. No se cuenta todo el caudal emitido como capturado.

\textbf{Modelo y ecuaciones.}

\[\fitmath{\dot m_c=\rho A(V-U),\quad F=\dot m_c(V-U),\quad P=FU,\quad |\boldsymbol v_{sal}|^2=U^2+(V-U)^2.}\]

\textbf{Desarrollo.}\begin{enumerate}

\item El flujo de masa a través de una frontera móvil usa la velocidad relativa $V-U$.

\item En laboratorio la componente longitudinal de salida es $U$; el cambio de momento da la fuerza sobre el carro.

\item Comparar potencia con la pérdida de energía cinética del caudal realmente interceptado.

\end{enumerate}

\textbf{Cálculo de cada magnitud.}

\begin{enumerate}

\item Obtener $\dot m_c$ con la relación del modelo y los valores anteriores.

\[\fitmath{\mathrm{md}=\rho\cdot A\cdot \left(V-U\right)}\]

\[\fitmath{\mathrm{md}=(1000)\cdot (0{,}0008)\cdot \left((24)-(6)\right)\approx 14{,}4\,\mathrm{kg/s}}\]

\item Obtener $F$ con la relación del modelo y los valores anteriores.

\[\fitmath{F=\mathrm{md}\cdot \left(V-U\right)}\]

\[\fitmath{F=(14{,}4)\cdot \left((24)-(6)\right)\approx 259{,}2\,\mathrm{N}}\]

\item Obtener $P$ con la relación del modelo y los valores anteriores.

\[\fitmath{\mathrm{power}=F\cdot U}\]

\[\fitmath{\mathrm{power}=(259{,}2)\cdot (6)\approx 1555{,}2\,\mathrm{W}}\]

\item Obtener $\Delta\dot E_c$ con la relación del modelo y los valores anteriores.

\[\fitmath{\mathrm{energy}=\frac{\mathrm{md}\cdot \left({V}^{2}-{U}^{2}-{\left(V-U\right)}^{2}\right)}{2}}\]

\[\fitmath{\mathrm{energy}=\frac{(14{,}4)\cdot \left({(24)}^{2}-{(6)}^{2}-{\left((24)-(6)\right)}^{2}\right)}{2}\approx 1555{,}2\,\mathrm{W}}\]

\end{enumerate}

\textbf{Resultados.}\begin{itemize}

\item $\dot m_c=14{,}4\,\mathrm{kg/s}$.

\item $F=259{,}2\,\mathrm{N}$.

\item $P=1555{,}2\,\mathrm{W}$.

\item $\Delta\dot E_c=1555{,}2\,\mathrm{W}$.

\end{itemize}

\textbf{Interpretación.} La velocidad relativa interviene dos veces en la fuerza: en masa capturada y en cambio de velocidad longitudinal.

\textbf{Comprobación.} Comprobar $\Delta\dot E_c=FU$. Si $U=0$, hay fuerza pero no potencia entregada al carro; si $U\to V$, desaparece la captura.

\textbf{Límites.} Un deflector que devuelve el chorro ciento ochenta grados tiene otro cambio de momento y otra fuerza.

\end{studiumworked}

\Needspace{.30\textheight}\section{MF-S4-4. Arranque impulsivo de una placa}

\begin{studiumworked}{MF-S4-4 · Intermedio}

\textbf{Enunciado.} Una placa infinita arranca instantáneamente con velocidad $U$. Determine velocidad en el líquido, esfuerzo en pared, transporte volumétrico por anchura y escala de penetración al tiempo dado.

\begin{longtable}{@{}p{.15\linewidth}p{.23\linewidth}p{.54\linewidth}@{}}
\toprule Símbolo & Valor SI & Significado\\\midrule\endhead
$U$ & $0{,}3\,\mathrm{m/s}$ & velocidad de placa\\
$\nu$ & $1\times10^{-6}\,\mathrm{m^2/s}$ & viscosidad cinemática\\
$\rho$ & $1000\,\mathrm{kg/m^3}$ & densidad\\
$y$ & $0{,}0012\,\mathrm{m}$ & distancia a placa\\
$t$ & $0{,}8\,\mathrm{s}$ & tiempo desde arranque\\
\bottomrule\end{longtable}

\textbf{Incógnitas.} Se calculan $u,\ \tau_{xy}(0),\ q,\ \delta_s$ y se interpreta el resultado bajo las hipótesis indicadas.

\textbf{Hipótesis.} Líquido newtoniano semiinfinito inicialmente en reposo, flujo paralelo sin gradiente de presión y $t>0$.

\textbf{Modelo y ecuaciones.}

\[\fitmath{u=U\,\operatorname{erfc}\!\left(\frac{y}{2\sqrt{\nu t}}\right),\quad \tau_{xy}(0)=-\frac{\mu U}{\sqrt{\pi\nu t}},\quad q=2U\sqrt{\nu t/\pi}.}\]

\textbf{Desarrollo.}\begin{enumerate}

\item La variable de semejanza es $y/(2\sqrt{\nu t})$. Aplicar las condiciones en placa, infinito e instante inicial para $y>0$.

\item Derivar la función error complementaria; el signo negativo corresponde a $\partial u/\partial y<0$.

\item Integrar el perfil en el semiespacio. La escala $2\sqrt{\nu t}$ no es un borde exacto del fluido en movimiento.

\end{enumerate}

\textbf{Cálculo de cada magnitud.}

\begin{enumerate}

\item Obtener $u$ con la relación del modelo y los valores anteriores.

\[\fitmath{u=U\cdot \operatorname{erfc}\left(\frac{y}{2\cdot \sqrt{\nu\cdot t}}\right)}\]

\[\fitmath{u=(0{,}3)\cdot \operatorname{erfc}\left(\frac{(0{,}0012)}{2\cdot \sqrt{(1\times10^{-6})\cdot (0{,}8)}}\right)\approx 0{,}1028345\,\mathrm{m/s}}\]

\item Obtener $\tau_{xy}(0)$ con la relación del modelo y los valores anteriores.

\[\fitmath{\tau=\frac{-\rho\cdot \nu\cdot U}{\sqrt{\pi\cdot \nu\cdot t}}}\]

\[\fitmath{\tau=\frac{-(1000)\cdot (1\times10^{-6})\cdot (0{,}3)}{\sqrt{\pi\cdot (1\times10^{-6})\cdot (0{,}8)}}\approx -0{,}1892349\,\mathrm{Pa}}\]

\item Obtener $q$ con la relación del modelo y los valores anteriores.

\[\fitmath{q=2\cdot U\cdot \sqrt{\frac{\nu\cdot t}{\pi}}}\]

\[\fitmath{q=2\cdot (0{,}3)\cdot \sqrt{\frac{(1\times10^{-6})\cdot (0{,}8)}{\pi}}\approx 0{,}0003027759\,\mathrm{m^{2}/s}}\]

\item Obtener $\delta_s$ con la relación del modelo y los valores anteriores.

\[\fitmath{\delta=2\cdot \sqrt{\nu\cdot t}}\]

\[\fitmath{\delta=2\cdot \sqrt{(1\times10^{-6})\cdot (0{,}8)}\approx 0{,}001788854\,\mathrm{m}}\]

\end{enumerate}

\textbf{Resultados.}\begin{itemize}

\item $u=0{,}102835\,\mathrm{m/s}$.

\item $\tau_{xy}(0)=-0{,}189235\,\mathrm{Pa}$.

\item $q=0{,}000302776\,\mathrm{m^2/s}$.

\item $\delta_s=0{,}00178885\,\mathrm{m}$.

\end{itemize}

\textbf{Interpretación.} El esfuerzo diverge al acercarse al instante del salto ideal de velocidad. Un arranque físico de duración finita elimina esa idealización.

\textbf{Comprobación.} Comprobar el residuo difusivo y el balance $\rho\,dq/dt=-\tau_{xy}(0)$.

\textbf{Límites.} No describe una placa finita, un arranque suave ni el efecto de otra pared cercana.

\end{studiumworked}

\Needspace{.30\textheight}\section{MF-S4-5. Convección térmica: escalas de Boussinesq}

\begin{studiumworked}{MF-S4-5 · Avanzado}

\textbf{Enunciado.} Con las propiedades académicas dadas, calcule difusividades, Prandtl, Grashof y Rayleigh. Compruebe la pequeña variación relativa de densidad y explique por qué Rayleigh por sí solo no fija un umbral universal de inestabilidad.

\begin{longtable}{@{}p{.15\linewidth}p{.23\linewidth}p{.54\linewidth}@{}}
\toprule Símbolo & Valor SI & Significado\\\midrule\endhead
$\rho$ & $1000\,\mathrm{kg/m^3}$ & densidad\\
$\mu$ & $0{,}001\,\mathrm{Pa\,s}$ & viscosidad\\
$k_T$ & $0{,}6\,\mathrm{W/(m\,K)}$ & conductividad\\
$c_p$ & $4200\,\mathrm{J/(kg\,K)}$ & calor específico\\
$\beta$ & $0{,}0002\,\mathrm{K^{-1}}$ & expansión térmica\\
$\Delta T$ & $4\,\mathrm{K}$ & diferencia térmica\\
$L$ & $0{,}012\,\mathrm{m}$ & longitud de referencia\\
\bottomrule\end{longtable}

\textbf{Incógnitas.} Se calculan $\nu,\ \kappa_T,\ \mathrm{Pr},\ \mathrm{Gr},\ \mathrm{Ra},\ \beta\Delta T$ y se interpreta el resultado bajo las hipótesis indicadas.

\textbf{Hipótesis.} Propiedades constantes de referencia, expansión lineal de densidad y gravedad uniforme. El enunciado define escalas, sin especificar una cavidad ni condiciones térmicas de contorno.

\textbf{Modelo y ecuaciones.}

\[\fitmath{\nu=\mu/\rho,\quad \kappa_T=k_T/(\rho c_p),\quad \mathrm{Pr}=\nu/\kappa_T,\quad\mathrm{Gr}=g\beta\Delta TL^3/\nu^2,\quad\mathrm{Ra}=\mathrm{Gr}\,\mathrm{Pr}.}\]

\textbf{Desarrollo.}\begin{enumerate}

\item Construir las dos difusividades con unidades de área por tiempo.

\item Formar los cocientes y verificar que son adimensionales.

\item La condición $\beta\Delta T\ll1$ apoya una parte de la aproximación Boussinesq. Un umbral crítico requiere geometría y condiciones de contorno concretas.

\end{enumerate}

\textbf{Cálculo de cada magnitud.}

\begin{enumerate}

\item Obtener $\nu$ con la relación del modelo y los valores anteriores.

\[\fitmath{\nu=\frac{\mu}{\rho}}\]

\[\fitmath{\nu=\frac{(0{,}001)}{(1000)}\approx 1\times10^{-6}\,\mathrm{m^{2}/s}}\]

\item Obtener $\kappa_T$ con la relación del modelo y los valores anteriores.

\[\fitmath{\kappa=\frac{\mathrm{kt}}{\rho\cdot \mathrm{cp}}}\]

\[\fitmath{\kappa=\frac{(0{,}6)}{(1000)\cdot (4200)}\approx 1{,}428571\times10^{-7}\,\mathrm{m^{2}/s}}\]

\item Obtener $\mathrm{Pr}$ con la relación del modelo y los valores anteriores.

\[\fitmath{\mathrm{Pr}=\frac{\nu}{\kappa}}\]

\[\fitmath{\mathrm{Pr}=\frac{(1\times10^{-6})}{(1{,}428571429\times10^{-7})}\approx 7\,}\]

\item Obtener $\mathrm{Gr}$ con la relación del modelo y los valores anteriores.

\[\fitmath{\mathrm{Gr}=\frac{g\cdot \beta\cdot \mathrm{dT}\cdot {L}^{3}}{{\nu}^{2}}}\]

\[\fitmath{\mathrm{Gr}=\frac{(9{,}81)\cdot (0{,}0002)\cdot (4)\cdot {(0{,}012)}^{3}}{{(1\times10^{-6})}^{2}}\approx 13561{,}34\,}\]

\item Obtener $\mathrm{Ra}$ con la relación del modelo y los valores anteriores.

\[\fitmath{\mathrm{Ra}=\mathrm{Gr}\cdot \mathrm{Pr}}\]

\[\fitmath{\mathrm{Ra}=(13561{,}344)\cdot (7)\approx 94929{,}41\,}\]

\item Obtener $\beta\Delta T$ con la relación del modelo y los valores anteriores.

\[\fitmath{\mathrm{variation}=\beta\cdot \mathrm{dT}}\]

\[\fitmath{\mathrm{variation}=(0{,}0002)\cdot (4)\approx 0{,}0008\,}\]

\end{enumerate}

\textbf{Resultados.}\begin{itemize}

\item $\nu=1\times10^{-6}\,\mathrm{m^2/s}$.

\item $\kappa_T=1{,}42857\times10^{-7}\,\mathrm{m^2/s}$.

\item $\mathrm{Pr}=7$.

\item $\mathrm{Gr}=13561{,}3$.

\item $\mathrm{Ra}=94929{,}4$.

\item $\beta\Delta T=0{,}0008$.

\end{itemize}

\textbf{Interpretación.} Rayleigh combina flotabilidad y dos procesos difusivos. No se aplica un valor crítico de una capa ideal a cualquier recinto o calentamiento.

\textbf{Comprobación.} Comprobar también $\mathrm{Ra}=g\beta\Delta TL^3/(\nu\kappa_T)$ y el límite sin diferencia térmica.

\textbf{Límites.} Densidad casi constante no asegura por sí sola todas las hipótesis de Boussinesq ni decide estabilidad sin el problema de contorno.

\end{studiumworked}

\studiumchaptercheck
